\documentclass[11pt]{article}

    \usepackage[breakable]{tcolorbox}
    \usepackage{parskip} % Stop auto-indenting (to mimic markdown behaviour)
    

    % Basic figure setup, for now with no caption control since it's done
    % automatically by Pandoc (which extracts ![](path) syntax from Markdown).
    \usepackage{graphicx}
    % Keep aspect ratio if custom image width or height is specified
    \setkeys{Gin}{keepaspectratio}
    % Maintain compatibility with old templates. Remove in nbconvert 6.0
    \let\Oldincludegraphics\includegraphics
    % Ensure that by default, figures have no caption (until we provide a
    % proper Figure object with a Caption API and a way to capture that
    % in the conversion process - todo).
    \usepackage{caption}
    \DeclareCaptionFormat{nocaption}{}
    \captionsetup{format=nocaption,aboveskip=0pt,belowskip=0pt}

    \usepackage{float}
    \floatplacement{figure}{H} % forces figures to be placed at the correct location
    \usepackage{xcolor} % Allow colors to be defined
    \usepackage{enumerate} % Needed for markdown enumerations to work
    \usepackage{geometry} % Used to adjust the document margins
    \usepackage{amsmath} % Equations
    \usepackage{amssymb} % Equations
    \usepackage{textcomp} % defines textquotesingle
    % Hack from http://tex.stackexchange.com/a/47451/13684:
    \AtBeginDocument{%
        \def\PYZsq{\textquotesingle}% Upright quotes in Pygmentized code
    }
    \usepackage{upquote} % Upright quotes for verbatim code
    \usepackage{eurosym} % defines \euro

    \usepackage{iftex}
    \ifPDFTeX
        \usepackage[T1]{fontenc}
        \IfFileExists{alphabeta.sty}{
              \usepackage{alphabeta}
          }{
              \usepackage[mathletters]{ucs}
              \usepackage[utf8x]{inputenc}
          }
    \else
        \usepackage{fontspec}
        \usepackage{unicode-math}
    \fi

    \usepackage{fancyvrb} % verbatim replacement that allows latex
    \usepackage{grffile} % extends the file name processing of package graphics
                         % to support a larger range
    \makeatletter % fix for old versions of grffile with XeLaTeX
    \@ifpackagelater{grffile}{2019/11/01}
    {
      % Do nothing on new versions
    }
    {
      \def\Gread@@xetex#1{%
        \IfFileExists{"\Gin@base".bb}%
        {\Gread@eps{\Gin@base.bb}}%
        {\Gread@@xetex@aux#1}%
      }
    }
    \makeatother
    \usepackage[Export]{adjustbox} % Used to constrain images to a maximum size
    \adjustboxset{max size={0.9\linewidth}{0.9\paperheight}}

    % The hyperref package gives us a pdf with properly built
    % internal navigation ('pdf bookmarks' for the table of contents,
    % internal cross-reference links, web links for URLs, etc.)
    \usepackage{hyperref}
    % The default LaTeX title has an obnoxious amount of whitespace. By default,
    % titling removes some of it. It also provides customization options.
    \usepackage{titling}
    \usepackage{longtable} % longtable support required by pandoc >1.10
    \usepackage{booktabs}  % table support for pandoc > 1.12.2
    \usepackage{array}     % table support for pandoc >= 2.11.3
    \usepackage{calc}      % table minipage width calculation for pandoc >= 2.11.1
    \usepackage[inline]{enumitem} % IRkernel/repr support (it uses the enumerate* environment)
    \usepackage[normalem]{ulem} % ulem is needed to support strikethroughs (\sout)
                                % normalem makes italics be italics, not underlines
    \usepackage{soul}      % strikethrough (\st) support for pandoc >= 3.0.0
    \usepackage{mathrsfs}
    

    
    % Colors for the hyperref package
    \definecolor{urlcolor}{rgb}{0,.145,.698}
    \definecolor{linkcolor}{rgb}{.71,0.21,0.01}
    \definecolor{citecolor}{rgb}{.12,.54,.11}

    % ANSI colors
    \definecolor{ansi-black}{HTML}{3E424D}
    \definecolor{ansi-black-intense}{HTML}{282C36}
    \definecolor{ansi-red}{HTML}{E75C58}
    \definecolor{ansi-red-intense}{HTML}{B22B31}
    \definecolor{ansi-green}{HTML}{00A250}
    \definecolor{ansi-green-intense}{HTML}{007427}
    \definecolor{ansi-yellow}{HTML}{DDB62B}
    \definecolor{ansi-yellow-intense}{HTML}{B27D12}
    \definecolor{ansi-blue}{HTML}{208FFB}
    \definecolor{ansi-blue-intense}{HTML}{0065CA}
    \definecolor{ansi-magenta}{HTML}{D160C4}
    \definecolor{ansi-magenta-intense}{HTML}{A03196}
    \definecolor{ansi-cyan}{HTML}{60C6C8}
    \definecolor{ansi-cyan-intense}{HTML}{258F8F}
    \definecolor{ansi-white}{HTML}{C5C1B4}
    \definecolor{ansi-white-intense}{HTML}{A1A6B2}
    \definecolor{ansi-default-inverse-fg}{HTML}{FFFFFF}
    \definecolor{ansi-default-inverse-bg}{HTML}{000000}

    % common color for the border for error outputs.
    \definecolor{outerrorbackground}{HTML}{FFDFDF}

    % commands and environments needed by pandoc snippets
    % extracted from the output of `pandoc -s`
    \providecommand{\tightlist}{%
      \setlength{\itemsep}{0pt}\setlength{\parskip}{0pt}}
    \DefineVerbatimEnvironment{Highlighting}{Verbatim}{commandchars=\\\{\}}
    % Add ',fontsize=\small' for more characters per line
    \newenvironment{Shaded}{}{}
    \newcommand{\KeywordTok}[1]{\textcolor[rgb]{0.00,0.44,0.13}{\textbf{{#1}}}}
    \newcommand{\DataTypeTok}[1]{\textcolor[rgb]{0.56,0.13,0.00}{{#1}}}
    \newcommand{\DecValTok}[1]{\textcolor[rgb]{0.25,0.63,0.44}{{#1}}}
    \newcommand{\BaseNTok}[1]{\textcolor[rgb]{0.25,0.63,0.44}{{#1}}}
    \newcommand{\FloatTok}[1]{\textcolor[rgb]{0.25,0.63,0.44}{{#1}}}
    \newcommand{\CharTok}[1]{\textcolor[rgb]{0.25,0.44,0.63}{{#1}}}
    \newcommand{\StringTok}[1]{\textcolor[rgb]{0.25,0.44,0.63}{{#1}}}
    \newcommand{\CommentTok}[1]{\textcolor[rgb]{0.38,0.63,0.69}{\textit{{#1}}}}
    \newcommand{\OtherTok}[1]{\textcolor[rgb]{0.00,0.44,0.13}{{#1}}}
    \newcommand{\AlertTok}[1]{\textcolor[rgb]{1.00,0.00,0.00}{\textbf{{#1}}}}
    \newcommand{\FunctionTok}[1]{\textcolor[rgb]{0.02,0.16,0.49}{{#1}}}
    \newcommand{\RegionMarkerTok}[1]{{#1}}
    \newcommand{\ErrorTok}[1]{\textcolor[rgb]{1.00,0.00,0.00}{\textbf{{#1}}}}
    \newcommand{\NormalTok}[1]{{#1}}

    % Additional commands for more recent versions of Pandoc
    \newcommand{\ConstantTok}[1]{\textcolor[rgb]{0.53,0.00,0.00}{{#1}}}
    \newcommand{\SpecialCharTok}[1]{\textcolor[rgb]{0.25,0.44,0.63}{{#1}}}
    \newcommand{\VerbatimStringTok}[1]{\textcolor[rgb]{0.25,0.44,0.63}{{#1}}}
    \newcommand{\SpecialStringTok}[1]{\textcolor[rgb]{0.73,0.40,0.53}{{#1}}}
    \newcommand{\ImportTok}[1]{{#1}}
    \newcommand{\DocumentationTok}[1]{\textcolor[rgb]{0.73,0.13,0.13}{\textit{{#1}}}}
    \newcommand{\AnnotationTok}[1]{\textcolor[rgb]{0.38,0.63,0.69}{\textbf{\textit{{#1}}}}}
    \newcommand{\CommentVarTok}[1]{\textcolor[rgb]{0.38,0.63,0.69}{\textbf{\textit{{#1}}}}}
    \newcommand{\VariableTok}[1]{\textcolor[rgb]{0.10,0.09,0.49}{{#1}}}
    \newcommand{\ControlFlowTok}[1]{\textcolor[rgb]{0.00,0.44,0.13}{\textbf{{#1}}}}
    \newcommand{\OperatorTok}[1]{\textcolor[rgb]{0.40,0.40,0.40}{{#1}}}
    \newcommand{\BuiltInTok}[1]{{#1}}
    \newcommand{\ExtensionTok}[1]{{#1}}
    \newcommand{\PreprocessorTok}[1]{\textcolor[rgb]{0.74,0.48,0.00}{{#1}}}
    \newcommand{\AttributeTok}[1]{\textcolor[rgb]{0.49,0.56,0.16}{{#1}}}
    \newcommand{\InformationTok}[1]{\textcolor[rgb]{0.38,0.63,0.69}{\textbf{\textit{{#1}}}}}
    \newcommand{\WarningTok}[1]{\textcolor[rgb]{0.38,0.63,0.69}{\textbf{\textit{{#1}}}}}
    \makeatletter
    \newsavebox\pandoc@box
    \newcommand*\pandocbounded[1]{%
      \sbox\pandoc@box{#1}%
      % scaling factors for width and height
      \Gscale@div\@tempa\textheight{\dimexpr\ht\pandoc@box+\dp\pandoc@box\relax}%
      \Gscale@div\@tempb\linewidth{\wd\pandoc@box}%
      % select the smaller of both
      \ifdim\@tempb\p@<\@tempa\p@
        \let\@tempa\@tempb
      \fi
      % scaling accordingly (\@tempa < 1)
      \ifdim\@tempa\p@<\p@
        \scalebox{\@tempa}{\usebox\pandoc@box}%
      % scaling not needed, use as it is
      \else
        \usebox{\pandoc@box}%
      \fi
    }
    \makeatother

    % Define a nice break command that doesn't care if a line doesn't already
    % exist.
    \def\br{\hspace*{\fill} \\* }
    % Math Jax compatibility definitions
    \def\gt{>}
    \def\lt{<}
    \let\Oldtex\TeX
    \let\Oldlatex\LaTeX
    \renewcommand{\TeX}{\textrm{\Oldtex}}
    \renewcommand{\LaTeX}{\textrm{\Oldlatex}}
    % Document parameters
    % Document title
    \title{stoichiometric\_justification}
    
    
    
    
    
    
    
% Pygments definitions
\makeatletter
\def\PY@reset{\let\PY@it=\relax \let\PY@bf=\relax%
    \let\PY@ul=\relax \let\PY@tc=\relax%
    \let\PY@bc=\relax \let\PY@ff=\relax}
\def\PY@tok#1{\csname PY@tok@#1\endcsname}
\def\PY@toks#1+{\ifx\relax#1\empty\else%
    \PY@tok{#1}\expandafter\PY@toks\fi}
\def\PY@do#1{\PY@bc{\PY@tc{\PY@ul{%
    \PY@it{\PY@bf{\PY@ff{#1}}}}}}}
\def\PY#1#2{\PY@reset\PY@toks#1+\relax+\PY@do{#2}}

\@namedef{PY@tok@w}{\def\PY@tc##1{\textcolor[rgb]{0.73,0.73,0.73}{##1}}}
\@namedef{PY@tok@c}{\let\PY@it=\textit\def\PY@tc##1{\textcolor[rgb]{0.24,0.48,0.48}{##1}}}
\@namedef{PY@tok@cp}{\def\PY@tc##1{\textcolor[rgb]{0.61,0.40,0.00}{##1}}}
\@namedef{PY@tok@k}{\let\PY@bf=\textbf\def\PY@tc##1{\textcolor[rgb]{0.00,0.50,0.00}{##1}}}
\@namedef{PY@tok@kp}{\def\PY@tc##1{\textcolor[rgb]{0.00,0.50,0.00}{##1}}}
\@namedef{PY@tok@kt}{\def\PY@tc##1{\textcolor[rgb]{0.69,0.00,0.25}{##1}}}
\@namedef{PY@tok@o}{\def\PY@tc##1{\textcolor[rgb]{0.40,0.40,0.40}{##1}}}
\@namedef{PY@tok@ow}{\let\PY@bf=\textbf\def\PY@tc##1{\textcolor[rgb]{0.67,0.13,1.00}{##1}}}
\@namedef{PY@tok@nb}{\def\PY@tc##1{\textcolor[rgb]{0.00,0.50,0.00}{##1}}}
\@namedef{PY@tok@nf}{\def\PY@tc##1{\textcolor[rgb]{0.00,0.00,1.00}{##1}}}
\@namedef{PY@tok@nc}{\let\PY@bf=\textbf\def\PY@tc##1{\textcolor[rgb]{0.00,0.00,1.00}{##1}}}
\@namedef{PY@tok@nn}{\let\PY@bf=\textbf\def\PY@tc##1{\textcolor[rgb]{0.00,0.00,1.00}{##1}}}
\@namedef{PY@tok@ne}{\let\PY@bf=\textbf\def\PY@tc##1{\textcolor[rgb]{0.80,0.25,0.22}{##1}}}
\@namedef{PY@tok@nv}{\def\PY@tc##1{\textcolor[rgb]{0.10,0.09,0.49}{##1}}}
\@namedef{PY@tok@no}{\def\PY@tc##1{\textcolor[rgb]{0.53,0.00,0.00}{##1}}}
\@namedef{PY@tok@nl}{\def\PY@tc##1{\textcolor[rgb]{0.46,0.46,0.00}{##1}}}
\@namedef{PY@tok@ni}{\let\PY@bf=\textbf\def\PY@tc##1{\textcolor[rgb]{0.44,0.44,0.44}{##1}}}
\@namedef{PY@tok@na}{\def\PY@tc##1{\textcolor[rgb]{0.41,0.47,0.13}{##1}}}
\@namedef{PY@tok@nt}{\let\PY@bf=\textbf\def\PY@tc##1{\textcolor[rgb]{0.00,0.50,0.00}{##1}}}
\@namedef{PY@tok@nd}{\def\PY@tc##1{\textcolor[rgb]{0.67,0.13,1.00}{##1}}}
\@namedef{PY@tok@s}{\def\PY@tc##1{\textcolor[rgb]{0.73,0.13,0.13}{##1}}}
\@namedef{PY@tok@sd}{\let\PY@it=\textit\def\PY@tc##1{\textcolor[rgb]{0.73,0.13,0.13}{##1}}}
\@namedef{PY@tok@si}{\let\PY@bf=\textbf\def\PY@tc##1{\textcolor[rgb]{0.64,0.35,0.47}{##1}}}
\@namedef{PY@tok@se}{\let\PY@bf=\textbf\def\PY@tc##1{\textcolor[rgb]{0.67,0.36,0.12}{##1}}}
\@namedef{PY@tok@sr}{\def\PY@tc##1{\textcolor[rgb]{0.64,0.35,0.47}{##1}}}
\@namedef{PY@tok@ss}{\def\PY@tc##1{\textcolor[rgb]{0.10,0.09,0.49}{##1}}}
\@namedef{PY@tok@sx}{\def\PY@tc##1{\textcolor[rgb]{0.00,0.50,0.00}{##1}}}
\@namedef{PY@tok@m}{\def\PY@tc##1{\textcolor[rgb]{0.40,0.40,0.40}{##1}}}
\@namedef{PY@tok@gh}{\let\PY@bf=\textbf\def\PY@tc##1{\textcolor[rgb]{0.00,0.00,0.50}{##1}}}
\@namedef{PY@tok@gu}{\let\PY@bf=\textbf\def\PY@tc##1{\textcolor[rgb]{0.50,0.00,0.50}{##1}}}
\@namedef{PY@tok@gd}{\def\PY@tc##1{\textcolor[rgb]{0.63,0.00,0.00}{##1}}}
\@namedef{PY@tok@gi}{\def\PY@tc##1{\textcolor[rgb]{0.00,0.52,0.00}{##1}}}
\@namedef{PY@tok@gr}{\def\PY@tc##1{\textcolor[rgb]{0.89,0.00,0.00}{##1}}}
\@namedef{PY@tok@ge}{\let\PY@it=\textit}
\@namedef{PY@tok@gs}{\let\PY@bf=\textbf}
\@namedef{PY@tok@ges}{\let\PY@bf=\textbf\let\PY@it=\textit}
\@namedef{PY@tok@gp}{\let\PY@bf=\textbf\def\PY@tc##1{\textcolor[rgb]{0.00,0.00,0.50}{##1}}}
\@namedef{PY@tok@go}{\def\PY@tc##1{\textcolor[rgb]{0.44,0.44,0.44}{##1}}}
\@namedef{PY@tok@gt}{\def\PY@tc##1{\textcolor[rgb]{0.00,0.27,0.87}{##1}}}
\@namedef{PY@tok@err}{\def\PY@bc##1{{\setlength{\fboxsep}{\string -\fboxrule}\fcolorbox[rgb]{1.00,0.00,0.00}{1,1,1}{\strut ##1}}}}
\@namedef{PY@tok@kc}{\let\PY@bf=\textbf\def\PY@tc##1{\textcolor[rgb]{0.00,0.50,0.00}{##1}}}
\@namedef{PY@tok@kd}{\let\PY@bf=\textbf\def\PY@tc##1{\textcolor[rgb]{0.00,0.50,0.00}{##1}}}
\@namedef{PY@tok@kn}{\let\PY@bf=\textbf\def\PY@tc##1{\textcolor[rgb]{0.00,0.50,0.00}{##1}}}
\@namedef{PY@tok@kr}{\let\PY@bf=\textbf\def\PY@tc##1{\textcolor[rgb]{0.00,0.50,0.00}{##1}}}
\@namedef{PY@tok@bp}{\def\PY@tc##1{\textcolor[rgb]{0.00,0.50,0.00}{##1}}}
\@namedef{PY@tok@fm}{\def\PY@tc##1{\textcolor[rgb]{0.00,0.00,1.00}{##1}}}
\@namedef{PY@tok@vc}{\def\PY@tc##1{\textcolor[rgb]{0.10,0.09,0.49}{##1}}}
\@namedef{PY@tok@vg}{\def\PY@tc##1{\textcolor[rgb]{0.10,0.09,0.49}{##1}}}
\@namedef{PY@tok@vi}{\def\PY@tc##1{\textcolor[rgb]{0.10,0.09,0.49}{##1}}}
\@namedef{PY@tok@vm}{\def\PY@tc##1{\textcolor[rgb]{0.10,0.09,0.49}{##1}}}
\@namedef{PY@tok@sa}{\def\PY@tc##1{\textcolor[rgb]{0.73,0.13,0.13}{##1}}}
\@namedef{PY@tok@sb}{\def\PY@tc##1{\textcolor[rgb]{0.73,0.13,0.13}{##1}}}
\@namedef{PY@tok@sc}{\def\PY@tc##1{\textcolor[rgb]{0.73,0.13,0.13}{##1}}}
\@namedef{PY@tok@dl}{\def\PY@tc##1{\textcolor[rgb]{0.73,0.13,0.13}{##1}}}
\@namedef{PY@tok@s2}{\def\PY@tc##1{\textcolor[rgb]{0.73,0.13,0.13}{##1}}}
\@namedef{PY@tok@sh}{\def\PY@tc##1{\textcolor[rgb]{0.73,0.13,0.13}{##1}}}
\@namedef{PY@tok@s1}{\def\PY@tc##1{\textcolor[rgb]{0.73,0.13,0.13}{##1}}}
\@namedef{PY@tok@mb}{\def\PY@tc##1{\textcolor[rgb]{0.40,0.40,0.40}{##1}}}
\@namedef{PY@tok@mf}{\def\PY@tc##1{\textcolor[rgb]{0.40,0.40,0.40}{##1}}}
\@namedef{PY@tok@mh}{\def\PY@tc##1{\textcolor[rgb]{0.40,0.40,0.40}{##1}}}
\@namedef{PY@tok@mi}{\def\PY@tc##1{\textcolor[rgb]{0.40,0.40,0.40}{##1}}}
\@namedef{PY@tok@il}{\def\PY@tc##1{\textcolor[rgb]{0.40,0.40,0.40}{##1}}}
\@namedef{PY@tok@mo}{\def\PY@tc##1{\textcolor[rgb]{0.40,0.40,0.40}{##1}}}
\@namedef{PY@tok@ch}{\let\PY@it=\textit\def\PY@tc##1{\textcolor[rgb]{0.24,0.48,0.48}{##1}}}
\@namedef{PY@tok@cm}{\let\PY@it=\textit\def\PY@tc##1{\textcolor[rgb]{0.24,0.48,0.48}{##1}}}
\@namedef{PY@tok@cpf}{\let\PY@it=\textit\def\PY@tc##1{\textcolor[rgb]{0.24,0.48,0.48}{##1}}}
\@namedef{PY@tok@c1}{\let\PY@it=\textit\def\PY@tc##1{\textcolor[rgb]{0.24,0.48,0.48}{##1}}}
\@namedef{PY@tok@cs}{\let\PY@it=\textit\def\PY@tc##1{\textcolor[rgb]{0.24,0.48,0.48}{##1}}}

\def\PYZbs{\char`\\}
\def\PYZus{\char`\_}
\def\PYZob{\char`\{}
\def\PYZcb{\char`\}}
\def\PYZca{\char`\^}
\def\PYZam{\char`\&}
\def\PYZlt{\char`\<}
\def\PYZgt{\char`\>}
\def\PYZsh{\char`\#}
\def\PYZpc{\char`\%}
\def\PYZdl{\char`\$}
\def\PYZhy{\char`\-}
\def\PYZsq{\char`\'}
\def\PYZdq{\char`\"}
\def\PYZti{\char`\~}
% for compatibility with earlier versions
\def\PYZat{@}
\def\PYZlb{[}
\def\PYZrb{]}
\makeatother


    % For linebreaks inside Verbatim environment from package fancyvrb.
    \makeatletter
        \newbox\Wrappedcontinuationbox
        \newbox\Wrappedvisiblespacebox
        \newcommand*\Wrappedvisiblespace {\textcolor{red}{\textvisiblespace}}
        \newcommand*\Wrappedcontinuationsymbol {\textcolor{red}{\llap{\tiny$\m@th\hookrightarrow$}}}
        \newcommand*\Wrappedcontinuationindent {3ex }
        \newcommand*\Wrappedafterbreak {\kern\Wrappedcontinuationindent\copy\Wrappedcontinuationbox}
        % Take advantage of the already applied Pygments mark-up to insert
        % potential linebreaks for TeX processing.
        %        {, <, #, %, $, ' and ": go to next line.
        %        _, }, ^, &, >, - and ~: stay at end of broken line.
        % Use of \textquotesingle for straight quote.
        \newcommand*\Wrappedbreaksatspecials {%
            \def\PYGZus{\discretionary{\char`\_}{\Wrappedafterbreak}{\char`\_}}%
            \def\PYGZob{\discretionary{}{\Wrappedafterbreak\char`\{}{\char`\{}}%
            \def\PYGZcb{\discretionary{\char`\}}{\Wrappedafterbreak}{\char`\}}}%
            \def\PYGZca{\discretionary{\char`\^}{\Wrappedafterbreak}{\char`\^}}%
            \def\PYGZam{\discretionary{\char`\&}{\Wrappedafterbreak}{\char`\&}}%
            \def\PYGZlt{\discretionary{}{\Wrappedafterbreak\char`\<}{\char`\<}}%
            \def\PYGZgt{\discretionary{\char`\>}{\Wrappedafterbreak}{\char`\>}}%
            \def\PYGZsh{\discretionary{}{\Wrappedafterbreak\char`\#}{\char`\#}}%
            \def\PYGZpc{\discretionary{}{\Wrappedafterbreak\char`\%}{\char`\%}}%
            \def\PYGZdl{\discretionary{}{\Wrappedafterbreak\char`\$}{\char`\$}}%
            \def\PYGZhy{\discretionary{\char`\-}{\Wrappedafterbreak}{\char`\-}}%
            \def\PYGZsq{\discretionary{}{\Wrappedafterbreak\textquotesingle}{\textquotesingle}}%
            \def\PYGZdq{\discretionary{}{\Wrappedafterbreak\char`\"}{\char`\"}}%
            \def\PYGZti{\discretionary{\char`\~}{\Wrappedafterbreak}{\char`\~}}%
        }
        % Some characters . , ; ? ! / are not pygmentized.
        % This macro makes them "active" and they will insert potential linebreaks
        \newcommand*\Wrappedbreaksatpunct {%
            \lccode`\~`\.\lowercase{\def~}{\discretionary{\hbox{\char`\.}}{\Wrappedafterbreak}{\hbox{\char`\.}}}%
            \lccode`\~`\,\lowercase{\def~}{\discretionary{\hbox{\char`\,}}{\Wrappedafterbreak}{\hbox{\char`\,}}}%
            \lccode`\~`\;\lowercase{\def~}{\discretionary{\hbox{\char`\;}}{\Wrappedafterbreak}{\hbox{\char`\;}}}%
            \lccode`\~`\:\lowercase{\def~}{\discretionary{\hbox{\char`\:}}{\Wrappedafterbreak}{\hbox{\char`\:}}}%
            \lccode`\~`\?\lowercase{\def~}{\discretionary{\hbox{\char`\?}}{\Wrappedafterbreak}{\hbox{\char`\?}}}%
            \lccode`\~`\!\lowercase{\def~}{\discretionary{\hbox{\char`\!}}{\Wrappedafterbreak}{\hbox{\char`\!}}}%
            \lccode`\~`\/\lowercase{\def~}{\discretionary{\hbox{\char`\/}}{\Wrappedafterbreak}{\hbox{\char`\/}}}%
            \catcode`\.\active
            \catcode`\,\active
            \catcode`\;\active
            \catcode`\:\active
            \catcode`\?\active
            \catcode`\!\active
            \catcode`\/\active
            \lccode`\~`\~
        }
    \makeatother

    \let\OriginalVerbatim=\Verbatim
    \makeatletter
    \renewcommand{\Verbatim}[1][1]{%
        %\parskip\z@skip
        \sbox\Wrappedcontinuationbox {\Wrappedcontinuationsymbol}%
        \sbox\Wrappedvisiblespacebox {\FV@SetupFont\Wrappedvisiblespace}%
        \def\FancyVerbFormatLine ##1{\hsize\linewidth
            \vtop{\raggedright\hyphenpenalty\z@\exhyphenpenalty\z@
                \doublehyphendemerits\z@\finalhyphendemerits\z@
                \strut ##1\strut}%
        }%
        % If the linebreak is at a space, the latter will be displayed as visible
        % space at end of first line, and a continuation symbol starts next line.
        % Stretch/shrink are however usually zero for typewriter font.
        \def\FV@Space {%
            \nobreak\hskip\z@ plus\fontdimen3\font minus\fontdimen4\font
            \discretionary{\copy\Wrappedvisiblespacebox}{\Wrappedafterbreak}
            {\kern\fontdimen2\font}%
        }%

        % Allow breaks at special characters using \PYG... macros.
        \Wrappedbreaksatspecials
        % Breaks at punctuation characters . , ; ? ! and / need catcode=\active
        \OriginalVerbatim[#1,codes*=\Wrappedbreaksatpunct]%
    }
    \makeatother

    % Exact colors from NB
    \definecolor{incolor}{HTML}{303F9F}
    \definecolor{outcolor}{HTML}{D84315}
    \definecolor{cellborder}{HTML}{CFCFCF}
    \definecolor{cellbackground}{HTML}{F7F7F7}

    % prompt
    \makeatletter
    \newcommand{\boxspacing}{\kern\kvtcb@left@rule\kern\kvtcb@boxsep}
    \makeatother
    \newcommand{\prompt}[4]{
        {\ttfamily\llap{{\color{#2}[#3]:\hspace{3pt}#4}}\vspace{-\baselineskip}}
    }
    

    
    % Prevent overflowing lines due to hard-to-break entities
    \sloppy
    % Setup hyperref package
    \hypersetup{
      breaklinks=true,  % so long urls are correctly broken across lines
      colorlinks=true,
      urlcolor=urlcolor,
      linkcolor=linkcolor,
      citecolor=citecolor,
      }
    % Slightly bigger margins than the latex defaults
    
    \geometry{verbose,tmargin=1in,bmargin=1in,lmargin=1in,rmargin=1in}
    
    

\begin{document}
    
    \maketitle
    
    

    
    \section{Stoichiometric Analysis of KNSB Composite Sugar
Propellant}\label{stoichiometric-analysis-of-knsb-composite-sugar-propellant}

\subsection{Background}\label{background}

KNSB (Potassium Nitrate / Sorbitol) is a castable solid rocket
propellant whose oxidiser--fuel ratio must be carefully chosen to
balance performance, safety, and manufacturability. To determine that
ratio rigorously we need a \textbf{balanced combustion equation}.

The reaction under study is:

\[
\text{KNO}_{3(s)} + \text{C}_6\text{H}_{14}\text{O}_{6(s)}
\;\rightarrow\;
\text{CO}_{2(g)} + \text{CO}_{(g)} + \text{H}_2\text{O}_{(g)} + \text{H}_{2(g)}
+ \text{N}_{2(g)} + \text{K}_2\text{CO}_{3(g)} + \text{KOH}_{(g)}
\]

We will 1. implement a \textbf{Python translation of the Mazza nullspace
algorithm} (Mazza \& Canuto, \emph{Fundamental Chemistry with MATLAB},
Elsevier 2022, Chapter 2), and 2. use it to balance a thermochemically
consistent, simplified form of the equation, 3. derive the ideal
stoichiometric mass ratio, and 4. justify why a slightly fuel-rich blend
is preferred in practice.

\emph{Ferric oxide (Fe₂O₃) acts as a catalyst and does not appear in the
equation.}

    \begin{center}\rule{0.5\linewidth}{0.5pt}\end{center}

\subsection{1. Python port of the Mazza nullspace
method}\label{python-port-of-the-mazza-nullspace-method}

\subsubsection{1.1 Theory (Chapter 2, Mazza \&
Canuto)}\label{theory-chapter-2-mazza-canuto}

Given \(n\) substances (reactants and products) containing \(m\)
elements, label the signed atom-count matrix
\(\mathbf{A} \in \mathbb{Z}^{m \times n}\):

\[
A_{jk} = \begin{cases}
+a_{jk} & \text{substance } k \text{ is a reactant} \\
-a_{jk} & \text{substance } k \text{ is a product}
\end{cases}
\]

where \(a_{jk}\) is the number of atoms of element \(j\) in substance
\(k\). Atom conservation then reads

\[
\mathbf{A}\,\mathbf{x} = \mathbf{0},
\qquad \mathbf{x} = [x_1, \ldots, x_n]^\top \in \mathbb{Z}_{>0}^n
\]

The solution lives in the \textbf{nullspace} of \(\mathbf{A}\).\\
A unique (up to scale) positive integer solution exists when
\(\dim(\ker\mathbf{A}) = 1\), i.e.~the nullspace is one-dimensional.
This is the \emph{minimum-balance stoichiometry}.

Mazza's algorithm (adapted from MATLAB to Python/SymPy): 1. Parse each
chemical formula as a symbolic expression. 2. Build \(\mathbf{A}\) by
reading atom counts from the symbolic tree. 3. Compute
\(\ker\mathbf{A}\) exactly with \texttt{Matrix.nullspace()}. 4. Convert
the rational nullspace vector to the minimum positive-integer form.

    \subsubsection{1.2 Environment}\label{environment}

    \begin{tcolorbox}[breakable, size=fbox, boxrule=1pt, pad at break*=1mm,colback=cellbackground, colframe=cellborder]
\prompt{In}{incolor}{1}{\boxspacing}
\begin{Verbatim}[commandchars=\\\{\}]
\PY{k+kn}{from}\PY{+w}{ }\PY{n+nn}{sympy}\PY{+w}{ }\PY{k+kn}{import} \PY{n}{symbols}\PY{p}{,} \PY{n}{Symbol}\PY{p}{,} \PY{n}{Mul}\PY{p}{,} \PY{n}{Pow}\PY{p}{,} \PY{n}{Matrix}\PY{p}{,} \PY{n}{Rational}\PY{p}{,} \PY{n}{lcm} \PY{k}{as} \PY{n}{sym\PYZus{}lcm}
\PY{k+kn}{from}\PY{+w}{ }\PY{n+nn}{math}\PY{+w}{ }\PY{k+kn}{import} \PY{n}{gcd}
\PY{k+kn}{from}\PY{+w}{ }\PY{n+nn}{functools}\PY{+w}{ }\PY{k+kn}{import} \PY{n}{reduce}
\PY{k+kn}{from}\PY{+w}{ }\PY{n+nn}{fractions}\PY{+w}{ }\PY{k+kn}{import} \PY{n}{Fraction}
\end{Verbatim}
\end{tcolorbox}

    \subsubsection{1.3 Formula parser}\label{formula-parser}

Translates the idea behind MATLAB's \texttt{symvar} / \texttt{children}
traversal: every symbolic expression is either a bare \texttt{Symbol}, a
\texttt{Pow} (element with exponent), or a \texttt{Mul} (product of the
previous two). We walk the expression tree recursively and collect
\texttt{\{element:\ count\}} pairs.

    \begin{tcolorbox}[breakable, size=fbox, boxrule=1pt, pad at break*=1mm,colback=cellbackground, colframe=cellborder]
\prompt{In}{incolor}{ }{\boxspacing}
\begin{Verbatim}[commandchars=\\\{\}]
\PY{k}{def}\PY{+w}{ }\PY{n+nf}{\PYZus{}parse\PYZus{}formula}\PY{p}{(}\PY{n}{expr}\PY{p}{,} \PY{n}{elem\PYZus{}set}\PY{p}{:} \PY{n+nb}{set}\PY{p}{)} \PY{o}{\PYZhy{}}\PY{o}{\PYZgt{}} \PY{n+nb}{dict}\PY{p}{:}
\PY{+w}{    }\PY{l+s+sd}{\PYZdq{}\PYZdq{}\PYZdq{}}
\PY{l+s+sd}{    Recursively parse a SymPy symbolic expression into an atom\PYZhy{}count dict.}

\PY{l+s+sd}{    Mirrors the MATLAB `children()` traversal in Mazza (2.4.3)}

\PY{l+s+sd}{    Parameters}
\PY{l+s+sd}{    \PYZhy{}\PYZhy{}\PYZhy{}\PYZhy{}\PYZhy{}\PYZhy{}\PYZhy{}\PYZhy{}\PYZhy{}\PYZhy{}}
\PY{l+s+sd}{    expr     : SymPy expression  (Symbol, Pow, or Mul)}
\PY{l+s+sd}{    elem\PYZus{}set : set of SymPy symbols  — the chemical elements in the reaction}

\PY{l+s+sd}{    Returns}
\PY{l+s+sd}{    \PYZhy{}\PYZhy{}\PYZhy{}\PYZhy{}\PYZhy{}\PYZhy{}\PYZhy{}}
\PY{l+s+sd}{    dict \PYZob{}symbol: int\PYZcb{}  atom counts for this formula}
\PY{l+s+sd}{    \PYZdq{}\PYZdq{}\PYZdq{}}
    \PY{n}{counts}\PY{p}{:} \PY{n+nb}{dict} \PY{o}{=} \PY{p}{\PYZob{}}\PY{p}{\PYZcb{}}

    \PY{k}{if} \PY{n+nb}{isinstance}\PY{p}{(}\PY{n}{expr}\PY{p}{,} \PY{n}{Symbol}\PY{p}{)}\PY{p}{:}              \PY{c+c1}{\PYZsh{} e.g.  O, K, H}
        \PY{k}{if} \PY{n}{expr} \PY{o+ow}{in} \PY{n}{elem\PYZus{}set}\PY{p}{:}
            \PY{n}{counts}\PY{p}{[}\PY{n}{expr}\PY{p}{]} \PY{o}{=} \PY{l+m+mi}{1}

    \PY{k}{elif} \PY{n+nb}{isinstance}\PY{p}{(}\PY{n}{expr}\PY{p}{,} \PY{n}{Pow}\PY{p}{)}\PY{p}{:}              \PY{c+c1}{\PYZsh{} e.g.  O**3, H**14}
        \PY{n}{base}\PY{p}{,} \PY{n}{exp} \PY{o}{=} \PY{n}{expr}\PY{o}{.}\PY{n}{args}
        \PY{k}{if} \PY{n}{base} \PY{o+ow}{in} \PY{n}{elem\PYZus{}set}\PY{p}{:}
            \PY{n}{counts}\PY{p}{[}\PY{n}{base}\PY{p}{]} \PY{o}{=} \PY{n+nb}{int}\PY{p}{(}\PY{n}{exp}\PY{p}{)}

    \PY{k}{elif} \PY{n+nb}{isinstance}\PY{p}{(}\PY{n}{expr}\PY{p}{,} \PY{n}{Mul}\PY{p}{)}\PY{p}{:}              \PY{c+c1}{\PYZsh{} e.g.  K*N*O**3, C**6*H**14*O**6}
        \PY{k}{for} \PY{n}{factor} \PY{o+ow}{in} \PY{n}{expr}\PY{o}{.}\PY{n}{args}\PY{p}{:}
            \PY{k}{for} \PY{n}{elem}\PY{p}{,} \PY{n}{n} \PY{o+ow}{in} \PY{n}{\PYZus{}parse\PYZus{}formula}\PY{p}{(}\PY{n}{factor}\PY{p}{,} \PY{n}{elem\PYZus{}set}\PY{p}{)}\PY{o}{.}\PY{n}{items}\PY{p}{(}\PY{p}{)}\PY{p}{:}
                \PY{n}{counts}\PY{p}{[}\PY{n}{elem}\PY{p}{]} \PY{o}{=} \PY{n}{counts}\PY{o}{.}\PY{n}{get}\PY{p}{(}\PY{n}{elem}\PY{p}{,} \PY{l+m+mi}{0}\PY{p}{)} \PY{o}{+} \PY{n}{n}

    \PY{k}{return} \PY{n}{counts}
\end{Verbatim}
\end{tcolorbox}

    \subsubsection{\texorpdfstring{1.4 \texttt{stoichiometry()} --- the core
function}{1.4 stoichiometry() --- the core function}}\label{stoichiometry-the-core-function}

Direct Python equivalent of the \texttt{stoichiometry.m} function in
Mazza §2.5.1. The main structural difference from the MATLAB version: -
\texttt{Matrix.nullspace()} returns exact rational vectors; no need for
a separate \texttt{rat()} call. - We use \texttt{functools.reduce} +
\texttt{math.gcd} for the LCM/GCD reduction.

    \begin{tcolorbox}[breakable, size=fbox, boxrule=1pt, pad at break*=1mm,colback=cellbackground, colframe=cellborder]
\prompt{In}{incolor}{3}{\boxspacing}
\begin{Verbatim}[commandchars=\\\{\}]
\PY{k}{def}\PY{+w}{ }\PY{n+nf}{stoichiometry}\PY{p}{(}
    \PY{n}{elements}\PY{p}{:} \PY{n+nb}{list}\PY{p}{,}
    \PY{n}{substances}\PY{p}{:} \PY{n+nb}{list}\PY{p}{,}
    \PY{n}{n\PYZus{}reactants}\PY{p}{:} \PY{n+nb}{int}\PY{p}{,}
    \PY{n}{label}\PY{p}{:} \PY{n+nb}{str} \PY{o}{=} \PY{l+s+s2}{\PYZdq{}}\PY{l+s+s2}{Reaction}\PY{l+s+s2}{\PYZdq{}}\PY{p}{,}
\PY{p}{)} \PY{o}{\PYZhy{}}\PY{o}{\PYZgt{}} \PY{n+nb}{list}\PY{p}{:}
\PY{+w}{    }\PY{l+s+sd}{\PYZdq{}\PYZdq{}\PYZdq{}}
\PY{l+s+sd}{    Minimum\PYZhy{}balance stoichiometry via the Mazza nullspace algorithm.}

\PY{l+s+sd}{    Port of `stoichiometry.m` (Mazza \PYZam{} Canuto, Chapter 2) from MATLAB to Python/SymPy.}

\PY{l+s+sd}{    Parameters}
\PY{l+s+sd}{    \PYZhy{}\PYZhy{}\PYZhy{}\PYZhy{}\PYZhy{}\PYZhy{}\PYZhy{}\PYZhy{}\PYZhy{}\PYZhy{}}
\PY{l+s+sd}{    elements    : list of SymPy symbols — chemical elements present (e.g. [K, N, O, C, H])}
\PY{l+s+sd}{    substances  : list of SymPy expressions — all substances, reactants first, then products}
\PY{l+s+sd}{                  written as symbolic products: K*N*O**3,  C**6*H**14*O**6, etc.}
\PY{l+s+sd}{    n\PYZus{}reactants : int — number of reactant substances (prefix of `substances`)}
\PY{l+s+sd}{    label       : str — human\PYZhy{}readable reaction name for the printed header}

\PY{l+s+sd}{    Returns}
\PY{l+s+sd}{    \PYZhy{}\PYZhy{}\PYZhy{}\PYZhy{}\PYZhy{}\PYZhy{}\PYZhy{}}
\PY{l+s+sd}{    int\PYZus{}coeffs : list[int] — minimum positive integer stoichiometric coefficients,}
\PY{l+s+sd}{                 in the same order as `substances` (reactants then products)}

\PY{l+s+sd}{    Raises}
\PY{l+s+sd}{    \PYZhy{}\PYZhy{}\PYZhy{}\PYZhy{}\PYZhy{}\PYZhy{}}
\PY{l+s+sd}{    ValueError  — if the nullspace is empty (no solution) or multi\PYZhy{}dimensional}
\PY{l+s+sd}{                  (underdetermined system; try removing linearly dependent species)}
\PY{l+s+sd}{    \PYZdq{}\PYZdq{}\PYZdq{}}
    \PY{n}{m} \PY{o}{=} \PY{n+nb}{len}\PY{p}{(}\PY{n}{elements}\PY{p}{)}          \PY{c+c1}{\PYZsh{} number of elements}
    \PY{n}{n} \PY{o}{=} \PY{n+nb}{len}\PY{p}{(}\PY{n}{substances}\PY{p}{)}        \PY{c+c1}{\PYZsh{} number of substances}
    \PY{n}{elem\PYZus{}set} \PY{o}{=} \PY{n+nb}{set}\PY{p}{(}\PY{n}{elements}\PY{p}{)}

    \PY{n+nb}{print}\PY{p}{(}\PY{l+s+sa}{f}\PY{l+s+s2}{\PYZdq{}}\PY{l+s+se}{\PYZbs{}n}\PY{l+s+si}{\PYZob{}}\PY{l+s+s1}{\PYZsq{}}\PY{l+s+s1}{─}\PY{l+s+s1}{\PYZsq{}}\PY{o}{*}\PY{l+m+mi}{55}\PY{l+s+si}{\PYZcb{}}\PY{l+s+s2}{\PYZdq{}}\PY{p}{)}
    \PY{n+nb}{print}\PY{p}{(}\PY{l+s+sa}{f}\PY{l+s+s2}{\PYZdq{}}\PY{l+s+s2}{  Minimum\PYZhy{}balance stoichiometry  —  }\PY{l+s+si}{\PYZob{}}\PY{n}{label}\PY{l+s+si}{\PYZcb{}}\PY{l+s+s2}{\PYZdq{}}\PY{p}{)}
    \PY{n+nb}{print}\PY{p}{(}\PY{l+s+sa}{f}\PY{l+s+s2}{\PYZdq{}}\PY{l+s+si}{\PYZob{}}\PY{l+s+s1}{\PYZsq{}}\PY{l+s+s1}{─}\PY{l+s+s1}{\PYZsq{}}\PY{o}{*}\PY{l+m+mi}{55}\PY{l+s+si}{\PYZcb{}}\PY{l+s+s2}{\PYZdq{}}\PY{p}{)}
    \PY{n+nb}{print}\PY{p}{(}\PY{l+s+sa}{f}\PY{l+s+s2}{\PYZdq{}}\PY{l+s+s2}{  Elements : }\PY{l+s+si}{\PYZob{}}\PY{n}{m}\PY{l+s+si}{\PYZcb{}}\PY{l+s+s2}{   Substances : }\PY{l+s+si}{\PYZob{}}\PY{n}{n}\PY{l+s+si}{\PYZcb{}}\PY{l+s+s2}{\PYZdq{}}\PY{p}{)}

    \PY{c+c1}{\PYZsh{} \PYZhy{}\PYZhy{}\PYZhy{}\PYZhy{}\PYZhy{}\PYZhy{}\PYZhy{}\PYZhy{}\PYZhy{}\PYZhy{}\PYZhy{}\PYZhy{}\PYZhy{}\PYZhy{}\PYZhy{}\PYZhy{}\PYZhy{}\PYZhy{}\PYZhy{}\PYZhy{}\PYZhy{}\PYZhy{}\PYZhy{}\PYZhy{}\PYZhy{}\PYZhy{}\PYZhy{}\PYZhy{}\PYZhy{}\PYZhy{}\PYZhy{}\PYZhy{}\PYZhy{}\PYZhy{}\PYZhy{}\PYZhy{}\PYZhy{}\PYZhy{}\PYZhy{}\PYZhy{}\PYZhy{}\PYZhy{}\PYZhy{}\PYZhy{}\PYZhy{}\PYZhy{}\PYZhy{}\PYZhy{}\PYZhy{}\PYZhy{}\PYZhy{}\PYZhy{}\PYZhy{}\PYZhy{}\PYZhy{}\PYZhy{}\PYZhy{}\PYZhy{}\PYZhy{}\PYZhy{}\PYZhy{}\PYZhy{}\PYZhy{}\PYZhy{}\PYZhy{}\PYZhy{} \PYZsh{}}
    \PY{c+c1}{\PYZsh{} Step 1 – build the signed balance matrix A  (m × n)}
    \PY{c+c1}{\PYZsh{} Mazza eq. (2.16):  A\PYZus{}\PYZob{}jk\PYZcb{} = +count for reactants, −count for products}
    \PY{c+c1}{\PYZsh{} \PYZhy{}\PYZhy{}\PYZhy{}\PYZhy{}\PYZhy{}\PYZhy{}\PYZhy{}\PYZhy{}\PYZhy{}\PYZhy{}\PYZhy{}\PYZhy{}\PYZhy{}\PYZhy{}\PYZhy{}\PYZhy{}\PYZhy{}\PYZhy{}\PYZhy{}\PYZhy{}\PYZhy{}\PYZhy{}\PYZhy{}\PYZhy{}\PYZhy{}\PYZhy{}\PYZhy{}\PYZhy{}\PYZhy{}\PYZhy{}\PYZhy{}\PYZhy{}\PYZhy{}\PYZhy{}\PYZhy{}\PYZhy{}\PYZhy{}\PYZhy{}\PYZhy{}\PYZhy{}\PYZhy{}\PYZhy{}\PYZhy{}\PYZhy{}\PYZhy{}\PYZhy{}\PYZhy{}\PYZhy{}\PYZhy{}\PYZhy{}\PYZhy{}\PYZhy{}\PYZhy{}\PYZhy{}\PYZhy{}\PYZhy{}\PYZhy{}\PYZhy{}\PYZhy{}\PYZhy{}\PYZhy{}\PYZhy{}\PYZhy{}\PYZhy{}\PYZhy{}\PYZhy{} \PYZsh{}}
    \PY{n}{A} \PY{o}{=} \PY{n}{Matrix}\PY{o}{.}\PY{n}{zeros}\PY{p}{(}\PY{n}{m}\PY{p}{,} \PY{n}{n}\PY{p}{)}

    \PY{k}{for} \PY{n}{k}\PY{p}{,} \PY{n}{formula} \PY{o+ow}{in} \PY{n+nb}{enumerate}\PY{p}{(}\PY{n}{substances}\PY{p}{)}\PY{p}{:}
        \PY{n}{sign} \PY{o}{=} \PY{l+m+mi}{1} \PY{k}{if} \PY{n}{k} \PY{o}{\PYZlt{}} \PY{n}{n\PYZus{}reactants} \PY{k}{else} \PY{o}{\PYZhy{}}\PY{l+m+mi}{1}
        \PY{n}{atom\PYZus{}counts} \PY{o}{=} \PY{n}{\PYZus{}parse\PYZus{}formula}\PY{p}{(}\PY{n}{formula}\PY{p}{,} \PY{n}{elem\PYZus{}set}\PY{p}{)}
        \PY{k}{for} \PY{n}{j}\PY{p}{,} \PY{n}{elem} \PY{o+ow}{in} \PY{n+nb}{enumerate}\PY{p}{(}\PY{n}{elements}\PY{p}{)}\PY{p}{:}
            \PY{n}{A}\PY{p}{[}\PY{n}{j}\PY{p}{,} \PY{n}{k}\PY{p}{]} \PY{o}{=} \PY{n}{sign} \PY{o}{*} \PY{n}{atom\PYZus{}counts}\PY{o}{.}\PY{n}{get}\PY{p}{(}\PY{n}{elem}\PY{p}{,} \PY{l+m+mi}{0}\PY{p}{)}

    \PY{c+c1}{\PYZsh{} \PYZhy{}\PYZhy{}\PYZhy{}\PYZhy{}\PYZhy{}\PYZhy{}\PYZhy{}\PYZhy{}\PYZhy{}\PYZhy{}\PYZhy{}\PYZhy{}\PYZhy{}\PYZhy{}\PYZhy{}\PYZhy{}\PYZhy{}\PYZhy{}\PYZhy{}\PYZhy{}\PYZhy{}\PYZhy{}\PYZhy{}\PYZhy{}\PYZhy{}\PYZhy{}\PYZhy{}\PYZhy{}\PYZhy{}\PYZhy{}\PYZhy{}\PYZhy{}\PYZhy{}\PYZhy{}\PYZhy{}\PYZhy{}\PYZhy{}\PYZhy{}\PYZhy{}\PYZhy{}\PYZhy{}\PYZhy{}\PYZhy{}\PYZhy{}\PYZhy{}\PYZhy{}\PYZhy{}\PYZhy{}\PYZhy{}\PYZhy{}\PYZhy{}\PYZhy{}\PYZhy{}\PYZhy{}\PYZhy{}\PYZhy{}\PYZhy{}\PYZhy{}\PYZhy{}\PYZhy{}\PYZhy{}\PYZhy{}\PYZhy{}\PYZhy{}\PYZhy{}\PYZhy{} \PYZsh{}}
    \PY{c+c1}{\PYZsh{} Step 2 – nullspace (exact, rational)}
    \PY{c+c1}{\PYZsh{} MATLAB: Z = null(A, \PYZsq{}r\PYZsq{})  →  SymPy: A.nullspace()}
    \PY{c+c1}{\PYZsh{} \PYZhy{}\PYZhy{}\PYZhy{}\PYZhy{}\PYZhy{}\PYZhy{}\PYZhy{}\PYZhy{}\PYZhy{}\PYZhy{}\PYZhy{}\PYZhy{}\PYZhy{}\PYZhy{}\PYZhy{}\PYZhy{}\PYZhy{}\PYZhy{}\PYZhy{}\PYZhy{}\PYZhy{}\PYZhy{}\PYZhy{}\PYZhy{}\PYZhy{}\PYZhy{}\PYZhy{}\PYZhy{}\PYZhy{}\PYZhy{}\PYZhy{}\PYZhy{}\PYZhy{}\PYZhy{}\PYZhy{}\PYZhy{}\PYZhy{}\PYZhy{}\PYZhy{}\PYZhy{}\PYZhy{}\PYZhy{}\PYZhy{}\PYZhy{}\PYZhy{}\PYZhy{}\PYZhy{}\PYZhy{}\PYZhy{}\PYZhy{}\PYZhy{}\PYZhy{}\PYZhy{}\PYZhy{}\PYZhy{}\PYZhy{}\PYZhy{}\PYZhy{}\PYZhy{}\PYZhy{}\PYZhy{}\PYZhy{}\PYZhy{}\PYZhy{}\PYZhy{}\PYZhy{} \PYZsh{}}
    \PY{n}{null\PYZus{}vecs} \PY{o}{=} \PY{n}{A}\PY{o}{.}\PY{n}{nullspace}\PY{p}{(}\PY{p}{)}
    \PY{n}{dim\PYZus{}null} \PY{o}{=} \PY{n+nb}{len}\PY{p}{(}\PY{n}{null\PYZus{}vecs}\PY{p}{)}

    \PY{k}{if} \PY{n}{dim\PYZus{}null} \PY{o}{==} \PY{l+m+mi}{0}\PY{p}{:}
        \PY{k}{raise} \PY{n+ne}{ValueError}\PY{p}{(}
            \PY{l+s+s2}{\PYZdq{}}\PY{l+s+s2}{Nullspace is empty — the system is inconsistent. }\PY{l+s+s2}{\PYZdq{}}
            \PY{l+s+s2}{\PYZdq{}}\PY{l+s+s2}{Check that all elements and species are correctly specified.}\PY{l+s+s2}{\PYZdq{}}
        \PY{p}{)}
    \PY{k}{if} \PY{n}{dim\PYZus{}null} \PY{o}{\PYZgt{}} \PY{l+m+mi}{1}\PY{p}{:}
        \PY{n+nb}{print}\PY{p}{(}
            \PY{l+s+sa}{f}\PY{l+s+s2}{\PYZdq{}}\PY{l+s+s2}{  ⚠  WARNING: nullspace dimension = }\PY{l+s+si}{\PYZob{}}\PY{n}{dim\PYZus{}null}\PY{l+s+si}{\PYZcb{}}\PY{l+s+s2}{ \PYZgt{} 1 — system is underdetermined.}\PY{l+s+se}{\PYZbs{}n}\PY{l+s+s2}{\PYZdq{}}
            \PY{l+s+sa}{f}\PY{l+s+s2}{\PYZdq{}}\PY{l+s+s2}{     The reaction as written has linearly dependent species and cannot}\PY{l+s+se}{\PYZbs{}n}\PY{l+s+s2}{\PYZdq{}}
            \PY{l+s+sa}{f}\PY{l+s+s2}{\PYZdq{}}\PY{l+s+s2}{     be given unique stoichiometric coefficients. Simplify the product set.}\PY{l+s+s2}{\PYZdq{}}
        \PY{p}{)}
        \PY{k}{return} \PY{k+kc}{None}

    \PY{n}{z} \PY{o}{=} \PY{n}{null\PYZus{}vecs}\PY{p}{[}\PY{l+m+mi}{0}\PY{p}{]}   \PY{c+c1}{\PYZsh{} column vector of Rational entries}

    \PY{c+c1}{\PYZsh{} \PYZhy{}\PYZhy{}\PYZhy{}\PYZhy{}\PYZhy{}\PYZhy{}\PYZhy{}\PYZhy{}\PYZhy{}\PYZhy{}\PYZhy{}\PYZhy{}\PYZhy{}\PYZhy{}\PYZhy{}\PYZhy{}\PYZhy{}\PYZhy{}\PYZhy{}\PYZhy{}\PYZhy{}\PYZhy{}\PYZhy{}\PYZhy{}\PYZhy{}\PYZhy{}\PYZhy{}\PYZhy{}\PYZhy{}\PYZhy{}\PYZhy{}\PYZhy{}\PYZhy{}\PYZhy{}\PYZhy{}\PYZhy{}\PYZhy{}\PYZhy{}\PYZhy{}\PYZhy{}\PYZhy{}\PYZhy{}\PYZhy{}\PYZhy{}\PYZhy{}\PYZhy{}\PYZhy{}\PYZhy{}\PYZhy{}\PYZhy{}\PYZhy{}\PYZhy{}\PYZhy{}\PYZhy{}\PYZhy{}\PYZhy{}\PYZhy{}\PYZhy{}\PYZhy{}\PYZhy{}\PYZhy{}\PYZhy{}\PYZhy{}\PYZhy{}\PYZhy{}\PYZhy{} \PYZsh{}}
    \PY{c+c1}{\PYZsh{} Step 3 – convert to minimum positive integers}
    \PY{c+c1}{\PYZsh{} MATLAB: [nZ,dZ]=rat(Z);  intZ=Z*lcm(sym(dZ))}
    \PY{c+c1}{\PYZsh{} \PYZhy{}\PYZhy{}\PYZhy{}\PYZhy{}\PYZhy{}\PYZhy{}\PYZhy{}\PYZhy{}\PYZhy{}\PYZhy{}\PYZhy{}\PYZhy{}\PYZhy{}\PYZhy{}\PYZhy{}\PYZhy{}\PYZhy{}\PYZhy{}\PYZhy{}\PYZhy{}\PYZhy{}\PYZhy{}\PYZhy{}\PYZhy{}\PYZhy{}\PYZhy{}\PYZhy{}\PYZhy{}\PYZhy{}\PYZhy{}\PYZhy{}\PYZhy{}\PYZhy{}\PYZhy{}\PYZhy{}\PYZhy{}\PYZhy{}\PYZhy{}\PYZhy{}\PYZhy{}\PYZhy{}\PYZhy{}\PYZhy{}\PYZhy{}\PYZhy{}\PYZhy{}\PYZhy{}\PYZhy{}\PYZhy{}\PYZhy{}\PYZhy{}\PYZhy{}\PYZhy{}\PYZhy{}\PYZhy{}\PYZhy{}\PYZhy{}\PYZhy{}\PYZhy{}\PYZhy{}\PYZhy{}\PYZhy{}\PYZhy{}\PYZhy{}\PYZhy{}\PYZhy{} \PYZsh{}}
    \PY{n}{fracs} \PY{o}{=} \PY{p}{[}\PY{n}{Fraction}\PY{p}{(}\PY{n+nb}{str}\PY{p}{(}\PY{n}{val}\PY{p}{)}\PY{p}{)} \PY{k}{for} \PY{n}{val} \PY{o+ow}{in} \PY{n}{z}\PY{p}{]}          \PY{c+c1}{\PYZsh{} exact fractions}
    \PY{n}{denoms} \PY{o}{=} \PY{p}{[}\PY{n}{f}\PY{o}{.}\PY{n}{denominator} \PY{k}{for} \PY{n}{f} \PY{o+ow}{in} \PY{n}{fracs}\PY{p}{]}

    \PY{k}{def}\PY{+w}{ }\PY{n+nf}{lcm2}\PY{p}{(}\PY{n}{a}\PY{p}{:} \PY{n+nb}{int}\PY{p}{,} \PY{n}{b}\PY{p}{:} \PY{n+nb}{int}\PY{p}{)} \PY{o}{\PYZhy{}}\PY{o}{\PYZgt{}} \PY{n+nb}{int}\PY{p}{:}
        \PY{k}{return} \PY{n}{a} \PY{o}{*} \PY{n}{b} \PY{o}{/}\PY{o}{/} \PY{n}{gcd}\PY{p}{(}\PY{n}{a}\PY{p}{,} \PY{n}{b}\PY{p}{)}

    \PY{n}{denom\PYZus{}lcm} \PY{o}{=} \PY{n}{reduce}\PY{p}{(}\PY{n}{lcm2}\PY{p}{,} \PY{n}{denoms}\PY{p}{)}                   \PY{c+c1}{\PYZsh{} LCM of denominators}
    \PY{n}{int\PYZus{}z} \PY{o}{=} \PY{p}{[}\PY{n+nb}{int}\PY{p}{(}\PY{n}{f} \PY{o}{*} \PY{n}{denom\PYZus{}lcm}\PY{p}{)} \PY{k}{for} \PY{n}{f} \PY{o+ow}{in} \PY{n}{fracs}\PY{p}{]}       \PY{c+c1}{\PYZsh{} scaled to integers}

    \PY{c+c1}{\PYZsh{} Ensure all positive}
    \PY{k}{if} \PY{n+nb}{any}\PY{p}{(}\PY{n}{v} \PY{o}{\PYZlt{}} \PY{l+m+mi}{0} \PY{k}{for} \PY{n}{v} \PY{o+ow}{in} \PY{n}{int\PYZus{}z}\PY{p}{)}\PY{p}{:}
        \PY{n}{int\PYZus{}z} \PY{o}{=} \PY{p}{[}\PY{o}{\PYZhy{}}\PY{n}{v} \PY{k}{for} \PY{n}{v} \PY{o+ow}{in} \PY{n}{int\PYZus{}z}\PY{p}{]}

    \PY{c+c1}{\PYZsh{} Divide by GCD to reach minimum norm}
    \PY{n}{g} \PY{o}{=} \PY{n}{reduce}\PY{p}{(}\PY{n}{gcd}\PY{p}{,} \PY{n}{int\PYZus{}z}\PY{p}{)}
    \PY{n}{int\PYZus{}z} \PY{o}{=} \PY{p}{[}\PY{n}{v} \PY{o}{/}\PY{o}{/} \PY{n}{g} \PY{k}{for} \PY{n}{v} \PY{o+ow}{in} \PY{n}{int\PYZus{}z}\PY{p}{]}

    \PY{c+c1}{\PYZsh{} \PYZhy{}\PYZhy{}\PYZhy{}\PYZhy{}\PYZhy{}\PYZhy{}\PYZhy{}\PYZhy{}\PYZhy{}\PYZhy{}\PYZhy{}\PYZhy{}\PYZhy{}\PYZhy{}\PYZhy{}\PYZhy{}\PYZhy{}\PYZhy{}\PYZhy{}\PYZhy{}\PYZhy{}\PYZhy{}\PYZhy{}\PYZhy{}\PYZhy{}\PYZhy{}\PYZhy{}\PYZhy{}\PYZhy{}\PYZhy{}\PYZhy{}\PYZhy{}\PYZhy{}\PYZhy{}\PYZhy{}\PYZhy{}\PYZhy{}\PYZhy{}\PYZhy{}\PYZhy{}\PYZhy{}\PYZhy{}\PYZhy{}\PYZhy{}\PYZhy{}\PYZhy{}\PYZhy{}\PYZhy{}\PYZhy{}\PYZhy{}\PYZhy{}\PYZhy{}\PYZhy{}\PYZhy{}\PYZhy{}\PYZhy{}\PYZhy{}\PYZhy{}\PYZhy{}\PYZhy{}\PYZhy{}\PYZhy{}\PYZhy{}\PYZhy{}\PYZhy{}\PYZhy{} \PYZsh{}}
    \PY{c+c1}{\PYZsh{} Step 4 – print results table}
    \PY{c+c1}{\PYZsh{} \PYZhy{}\PYZhy{}\PYZhy{}\PYZhy{}\PYZhy{}\PYZhy{}\PYZhy{}\PYZhy{}\PYZhy{}\PYZhy{}\PYZhy{}\PYZhy{}\PYZhy{}\PYZhy{}\PYZhy{}\PYZhy{}\PYZhy{}\PYZhy{}\PYZhy{}\PYZhy{}\PYZhy{}\PYZhy{}\PYZhy{}\PYZhy{}\PYZhy{}\PYZhy{}\PYZhy{}\PYZhy{}\PYZhy{}\PYZhy{}\PYZhy{}\PYZhy{}\PYZhy{}\PYZhy{}\PYZhy{}\PYZhy{}\PYZhy{}\PYZhy{}\PYZhy{}\PYZhy{}\PYZhy{}\PYZhy{}\PYZhy{}\PYZhy{}\PYZhy{}\PYZhy{}\PYZhy{}\PYZhy{}\PYZhy{}\PYZhy{}\PYZhy{}\PYZhy{}\PYZhy{}\PYZhy{}\PYZhy{}\PYZhy{}\PYZhy{}\PYZhy{}\PYZhy{}\PYZhy{}\PYZhy{}\PYZhy{}\PYZhy{}\PYZhy{}\PYZhy{}\PYZhy{} \PYZsh{}}
    \PY{n+nb}{print}\PY{p}{(}\PY{l+s+sa}{f}\PY{l+s+s2}{\PYZdq{}}\PY{l+s+se}{\PYZbs{}n}\PY{l+s+s2}{  }\PY{l+s+si}{\PYZob{}}\PY{l+s+s1}{\PYZsq{}}\PY{l+s+s1}{Substance}\PY{l+s+s1}{\PYZsq{}}\PY{l+s+si}{:}\PY{l+s+s2}{\PYZlt{}22}\PY{l+s+si}{\PYZcb{}}\PY{l+s+s2}{ }\PY{l+s+si}{\PYZob{}}\PY{l+s+s1}{\PYZsq{}}\PY{l+s+s1}{Coeff}\PY{l+s+s1}{\PYZsq{}}\PY{l+s+si}{:}\PY{l+s+s2}{\PYZgt{}6}\PY{l+s+si}{\PYZcb{}}\PY{l+s+s2}{  }\PY{l+s+si}{\PYZob{}}\PY{l+s+s1}{\PYZsq{}}\PY{l+s+s1}{Type}\PY{l+s+s1}{\PYZsq{}}\PY{l+s+si}{\PYZcb{}}\PY{l+s+s2}{\PYZdq{}}\PY{p}{)}
    \PY{n+nb}{print}\PY{p}{(}\PY{l+s+sa}{f}\PY{l+s+s2}{\PYZdq{}}\PY{l+s+s2}{  }\PY{l+s+si}{\PYZob{}}\PY{l+s+s1}{\PYZsq{}}\PY{l+s+s1}{─}\PY{l+s+s1}{\PYZsq{}}\PY{o}{*}\PY{l+m+mi}{22}\PY{l+s+si}{\PYZcb{}}\PY{l+s+s2}{ }\PY{l+s+si}{\PYZob{}}\PY{l+s+s1}{\PYZsq{}}\PY{l+s+s1}{─}\PY{l+s+s1}{\PYZsq{}}\PY{o}{*}\PY{l+m+mi}{6}\PY{l+s+si}{\PYZcb{}}\PY{l+s+s2}{  }\PY{l+s+si}{\PYZob{}}\PY{l+s+s1}{\PYZsq{}}\PY{l+s+s1}{─}\PY{l+s+s1}{\PYZsq{}}\PY{o}{*}\PY{l+m+mi}{8}\PY{l+s+si}{\PYZcb{}}\PY{l+s+s2}{\PYZdq{}}\PY{p}{)}
    \PY{k}{for} \PY{n}{k}\PY{p}{,} \PY{p}{(}\PY{n}{formula}\PY{p}{,} \PY{n}{coeff}\PY{p}{)} \PY{o+ow}{in} \PY{n+nb}{enumerate}\PY{p}{(}\PY{n+nb}{zip}\PY{p}{(}\PY{n}{substances}\PY{p}{,} \PY{n}{int\PYZus{}z}\PY{p}{)}\PY{p}{)}\PY{p}{:}
        \PY{n}{role} \PY{o}{=} \PY{l+s+s2}{\PYZdq{}}\PY{l+s+s2}{Reactant}\PY{l+s+s2}{\PYZdq{}} \PY{k}{if} \PY{n}{k} \PY{o}{\PYZlt{}} \PY{n}{n\PYZus{}reactants} \PY{k}{else} \PY{l+s+s2}{\PYZdq{}}\PY{l+s+s2}{Product}\PY{l+s+s2}{\PYZdq{}}
        \PY{n+nb}{print}\PY{p}{(}\PY{l+s+sa}{f}\PY{l+s+s2}{\PYZdq{}}\PY{l+s+s2}{  }\PY{l+s+si}{\PYZob{}}\PY{n+nb}{str}\PY{p}{(}\PY{n}{formula}\PY{p}{)}\PY{l+s+si}{:}\PY{l+s+s2}{\PYZlt{}22}\PY{l+s+si}{\PYZcb{}}\PY{l+s+s2}{ }\PY{l+s+si}{\PYZob{}}\PY{n}{coeff}\PY{l+s+si}{:}\PY{l+s+s2}{\PYZgt{}6}\PY{l+s+si}{\PYZcb{}}\PY{l+s+s2}{  }\PY{l+s+si}{\PYZob{}}\PY{n}{role}\PY{l+s+si}{\PYZcb{}}\PY{l+s+s2}{\PYZdq{}}\PY{p}{)}

    \PY{k}{return} \PY{n}{int\PYZus{}z}
\end{Verbatim}
\end{tcolorbox}

    \begin{center}\rule{0.5\linewidth}{0.5pt}\end{center}

\subsection{2. Full KNSB reaction --- infeasibility
check}\label{full-knsb-reaction-infeasibility-check}

The commonly cited full product set is:

\[
a\,\text{KNO}_3 + b\,\text{C}_6\text{H}_{14}\text{O}_6
\;\rightarrow\;
c\,\text{CO}_2 + d\,\text{CO} + e\,\text{H}_2\text{O} + f\,\text{H}_2
+ g\,\text{N}_2 + h\,\text{K}_2\text{CO}_3 + i\,\text{KOH}
\]

We have 5 elements and 9 species, so
\(\dim(\ker\mathbf{A}) = 9 - \text{rank}(\mathbf{A}) \geq 9 - 5 = 4\). A
unique balanced equation requires the nullspace to be 1-dimensional,
which is impossible here. The \texttt{stoichiometry()} call below
confirms this.

    \begin{tcolorbox}[breakable, size=fbox, boxrule=1pt, pad at break*=1mm,colback=cellbackground, colframe=cellborder]
\prompt{In}{incolor}{4}{\boxspacing}
\begin{Verbatim}[commandchars=\\\{\}]
\PY{c+c1}{\PYZsh{} ── Element symbols  (same convention as the MATLAB notebook) ──────────────}
\PY{n}{K}\PY{p}{,} \PY{n}{N}\PY{p}{,} \PY{n}{O}\PY{p}{,} \PY{n}{C}\PY{p}{,} \PY{n}{H} \PY{o}{=} \PY{n}{symbols}\PY{p}{(}\PY{l+s+s1}{\PYZsq{}}\PY{l+s+s1}{K N O C H}\PY{l+s+s1}{\PYZsq{}}\PY{p}{)}
\PY{n}{elements} \PY{o}{=} \PY{p}{[}\PY{n}{K}\PY{p}{,} \PY{n}{N}\PY{p}{,} \PY{n}{O}\PY{p}{,} \PY{n}{C}\PY{p}{,} \PY{n}{H}\PY{p}{]}

\PY{c+c1}{\PYZsh{} ── Substances: reactants then products ───────────────────────────────────}
\PY{c+c1}{\PYZsh{} KNO3           C6H14O6}
\PY{n}{reactants\PYZus{}full} \PY{o}{=} \PY{p}{[}\PY{n}{K}\PY{o}{*}\PY{n}{N}\PY{o}{*}\PY{n}{O}\PY{o}{*}\PY{o}{*}\PY{l+m+mi}{3}\PY{p}{,}  \PY{n}{C}\PY{o}{*}\PY{o}{*}\PY{l+m+mi}{6}\PY{o}{*}\PY{n}{H}\PY{o}{*}\PY{o}{*}\PY{l+m+mi}{14}\PY{o}{*}\PY{n}{O}\PY{o}{*}\PY{o}{*}\PY{l+m+mi}{6}\PY{p}{]}

\PY{c+c1}{\PYZsh{} CO2      CO     H2O      H2      N2      K2CO3          KOH}
\PY{n}{products\PYZus{}full} \PY{o}{=} \PY{p}{[}
    \PY{n}{C}\PY{o}{*}\PY{n}{O}\PY{o}{*}\PY{o}{*}\PY{l+m+mi}{2}\PY{p}{,} \PY{n}{C}\PY{o}{*}\PY{n}{O}\PY{p}{,} \PY{n}{H}\PY{o}{*}\PY{o}{*}\PY{l+m+mi}{2}\PY{o}{*}\PY{n}{O}\PY{p}{,} \PY{n}{H}\PY{o}{*}\PY{o}{*}\PY{l+m+mi}{2}\PY{p}{,} \PY{n}{N}\PY{o}{*}\PY{o}{*}\PY{l+m+mi}{2}\PY{p}{,}
    \PY{n}{K}\PY{o}{*}\PY{o}{*}\PY{l+m+mi}{2}\PY{o}{*}\PY{n}{C}\PY{o}{*}\PY{n}{O}\PY{o}{*}\PY{o}{*}\PY{l+m+mi}{3}\PY{p}{,}
    \PY{n}{K}\PY{o}{*}\PY{n}{O}\PY{o}{*}\PY{n}{H}
\PY{p}{]}

\PY{n}{result\PYZus{}full} \PY{o}{=} \PY{n}{stoichiometry}\PY{p}{(}
    \PY{n}{elements}\PY{p}{,}
    \PY{n}{reactants\PYZus{}full} \PY{o}{+} \PY{n}{products\PYZus{}full}\PY{p}{,}
    \PY{n}{n\PYZus{}reactants}\PY{o}{=}\PY{l+m+mi}{2}\PY{p}{,}
    \PY{n}{label}\PY{o}{=}\PY{l+s+s2}{\PYZdq{}}\PY{l+s+s2}{KNSB — full product set (9 species)}\PY{l+s+s2}{\PYZdq{}}
\PY{p}{)}
\end{Verbatim}
\end{tcolorbox}

    \begin{Verbatim}[commandchars=\\\{\}]

───────────────────────────────────────────────────────
  Minimum-balance stoichiometry  —  KNSB — full product set (9 species)
───────────────────────────────────────────────────────
  Elements : 5   Substances : 9
  ⚠  WARNING: nullspace dimension = 4 > 1 — system is underdetermined.
     The reaction as written has linearly dependent species and cannot
     be given unique stoichiometric coefficients. Simplify the product set.
    \end{Verbatim}

    \begin{quote}
\textbf{Interpretation}: The nullspace dimension equals the number of
free parameters in the product distribution. With 9 species and only 5
element constraints the system has a 4-dimensional nullspace --- no
unique balanced equation exists. We must fix the product slate to obtain
a physically meaningful result.
\end{quote}

    \begin{center}\rule{0.5\linewidth}{0.5pt}\end{center}

\subsection{3. Simplified KNSB reaction --- unique balanced
equation}\label{simplified-knsb-reaction-unique-balanced-equation}

Thermochemical analysis (and experimental validation by Nakka 2025)
shows that for near-stoichiometric KNSB at moderate chamber pressure,
the dominant products are \(\text{CO}_2\), \(\text{H}_2\text{O}\),
\(\text{N}_2\), and \(\text{K}_2\text{CO}_3\). Dropping the minority
species CO, H₂, and KOH reduces the system to 6 substances, giving a
unique (1-D nullspace) solution:

\[
a\,\text{KNO}_3 + b\,\text{C}_6\text{H}_{14}\text{O}_6
\;\rightarrow\;
c\,\text{CO}_2 + d\,\text{H}_2\text{O} + e\,\text{N}_2 + f\,\text{K}_2\text{CO}_3
\]

    \begin{tcolorbox}[breakable, size=fbox, boxrule=1pt, pad at break*=1mm,colback=cellbackground, colframe=cellborder]
\prompt{In}{incolor}{5}{\boxspacing}
\begin{Verbatim}[commandchars=\\\{\}]
\PY{c+c1}{\PYZsh{} ── Simplified product set ─────────────────────────────────────────────────}
\PY{c+c1}{\PYZsh{}        CO2      H2O       N2      K2CO3}
\PY{n}{products\PYZus{}simplified} \PY{o}{=} \PY{p}{[}\PY{n}{C}\PY{o}{*}\PY{n}{O}\PY{o}{*}\PY{o}{*}\PY{l+m+mi}{2}\PY{p}{,} \PY{n}{H}\PY{o}{*}\PY{o}{*}\PY{l+m+mi}{2}\PY{o}{*}\PY{n}{O}\PY{p}{,} \PY{n}{N}\PY{o}{*}\PY{o}{*}\PY{l+m+mi}{2}\PY{p}{,} \PY{n}{K}\PY{o}{*}\PY{o}{*}\PY{l+m+mi}{2}\PY{o}{*}\PY{n}{C}\PY{o}{*}\PY{n}{O}\PY{o}{*}\PY{o}{*}\PY{l+m+mi}{3}\PY{p}{]}

\PY{n}{int\PYZus{}coeffs} \PY{o}{=} \PY{n}{stoichiometry}\PY{p}{(}
    \PY{n}{elements}\PY{p}{,}
    \PY{n}{reactants\PYZus{}full} \PY{o}{+} \PY{n}{products\PYZus{}simplified}\PY{p}{,}
    \PY{n}{n\PYZus{}reactants}\PY{o}{=}\PY{l+m+mi}{2}\PY{p}{,}
    \PY{n}{label}\PY{o}{=}\PY{l+s+s2}{\PYZdq{}}\PY{l+s+s2}{KNSB — simplified product set (6 species)}\PY{l+s+s2}{\PYZdq{}}
\PY{p}{)}

\PY{c+c1}{\PYZsh{} Unpack results for later use}
\PY{n}{a\PYZus{}KNO3}\PY{p}{,} \PY{n}{b\PYZus{}C6H14O6}\PY{p}{,} \PY{n}{c\PYZus{}CO2}\PY{p}{,} \PY{n}{d\PYZus{}H2O}\PY{p}{,} \PY{n}{e\PYZus{}N2}\PY{p}{,} \PY{n}{f\PYZus{}K2CO3} \PY{o}{=} \PY{n}{int\PYZus{}coeffs}
\end{Verbatim}
\end{tcolorbox}

    \begin{Verbatim}[commandchars=\\\{\}]

───────────────────────────────────────────────────────
  Minimum-balance stoichiometry  —  KNSB — simplified product set (6 species)
───────────────────────────────────────────────────────
  Elements : 5   Substances : 6

  Substance               Coeff  Type
  ────────────────────── ──────  ────────
  K*N*O**3                   26  Reactant
  C**6*H**14*O**6             5  Reactant
  C*O**2                     17  Product
  H**2*O                     35  Product
  N**2                       13  Product
  C*K**2*O**3                13  Product
    \end{Verbatim}

    \begin{tcolorbox}[breakable, size=fbox, boxrule=1pt, pad at break*=1mm,colback=cellbackground, colframe=cellborder]
\prompt{In}{incolor}{6}{\boxspacing}
\begin{Verbatim}[commandchars=\\\{\}]
\PY{c+c1}{\PYZsh{} Pretty\PYZhy{}print the balanced equation}
\PY{n+nb}{print}\PY{p}{(}\PY{l+s+s2}{\PYZdq{}}\PY{l+s+se}{\PYZbs{}n}\PY{l+s+s2}{Balanced combustion equation:}\PY{l+s+s2}{\PYZdq{}}\PY{p}{)}
\PY{n+nb}{print}\PY{p}{(}
    \PY{l+s+sa}{f}\PY{l+s+s2}{\PYZdq{}}\PY{l+s+s2}{  }\PY{l+s+si}{\PYZob{}}\PY{n}{a\PYZus{}KNO3}\PY{l+s+si}{\PYZcb{}}\PY{l+s+s2}{ KNO₃  +  }\PY{l+s+si}{\PYZob{}}\PY{n}{b\PYZus{}C6H14O6}\PY{l+s+si}{\PYZcb{}}\PY{l+s+s2}{ C₆H₁₄O₆}\PY{l+s+s2}{\PYZdq{}}
    \PY{l+s+sa}{f}\PY{l+s+s2}{\PYZdq{}}\PY{l+s+s2}{  →  }\PY{l+s+si}{\PYZob{}}\PY{n}{c\PYZus{}CO2}\PY{l+s+si}{\PYZcb{}}\PY{l+s+s2}{ CO₂  +  }\PY{l+s+si}{\PYZob{}}\PY{n}{d\PYZus{}H2O}\PY{l+s+si}{\PYZcb{}}\PY{l+s+s2}{ H₂O}\PY{l+s+s2}{\PYZdq{}}
    \PY{l+s+sa}{f}\PY{l+s+s2}{\PYZdq{}}\PY{l+s+s2}{  +  }\PY{l+s+si}{\PYZob{}}\PY{n}{e\PYZus{}N2}\PY{l+s+si}{\PYZcb{}}\PY{l+s+s2}{ N₂  +  }\PY{l+s+si}{\PYZob{}}\PY{n}{f\PYZus{}K2CO3}\PY{l+s+si}{\PYZcb{}}\PY{l+s+s2}{ K₂CO₃}\PY{l+s+s2}{\PYZdq{}}
\PY{p}{)}

\PY{c+c1}{\PYZsh{} Verify atom conservation}
\PY{n+nb}{print}\PY{p}{(}\PY{l+s+s2}{\PYZdq{}}\PY{l+s+se}{\PYZbs{}n}\PY{l+s+s2}{Atom\PYZhy{}conservation check:}\PY{l+s+s2}{\PYZdq{}}\PY{p}{)}
\PY{n}{atoms\PYZus{}left}  \PY{o}{=} \PY{p}{\PYZob{}}\PY{l+s+s2}{\PYZdq{}}\PY{l+s+s2}{K}\PY{l+s+s2}{\PYZdq{}}\PY{p}{:} \PY{n}{a\PYZus{}KNO3}\PY{p}{,} \PY{l+s+s2}{\PYZdq{}}\PY{l+s+s2}{N}\PY{l+s+s2}{\PYZdq{}}\PY{p}{:} \PY{n}{a\PYZus{}KNO3}\PY{p}{,}
               \PY{l+s+s2}{\PYZdq{}}\PY{l+s+s2}{O}\PY{l+s+s2}{\PYZdq{}}\PY{p}{:} \PY{l+m+mi}{3}\PY{o}{*}\PY{n}{a\PYZus{}KNO3} \PY{o}{+} \PY{l+m+mi}{6}\PY{o}{*}\PY{n}{b\PYZus{}C6H14O6}\PY{p}{,}
               \PY{l+s+s2}{\PYZdq{}}\PY{l+s+s2}{C}\PY{l+s+s2}{\PYZdq{}}\PY{p}{:} \PY{l+m+mi}{6}\PY{o}{*}\PY{n}{b\PYZus{}C6H14O6}\PY{p}{,} \PY{l+s+s2}{\PYZdq{}}\PY{l+s+s2}{H}\PY{l+s+s2}{\PYZdq{}}\PY{p}{:} \PY{l+m+mi}{14}\PY{o}{*}\PY{n}{b\PYZus{}C6H14O6}\PY{p}{\PYZcb{}}
\PY{n}{atoms\PYZus{}right} \PY{o}{=} \PY{p}{\PYZob{}}\PY{l+s+s2}{\PYZdq{}}\PY{l+s+s2}{K}\PY{l+s+s2}{\PYZdq{}}\PY{p}{:} \PY{l+m+mi}{2}\PY{o}{*}\PY{n}{f\PYZus{}K2CO3}\PY{p}{,} \PY{l+s+s2}{\PYZdq{}}\PY{l+s+s2}{N}\PY{l+s+s2}{\PYZdq{}}\PY{p}{:} \PY{l+m+mi}{2}\PY{o}{*}\PY{n}{e\PYZus{}N2}\PY{p}{,}
               \PY{l+s+s2}{\PYZdq{}}\PY{l+s+s2}{O}\PY{l+s+s2}{\PYZdq{}}\PY{p}{:} \PY{l+m+mi}{2}\PY{o}{*}\PY{n}{c\PYZus{}CO2} \PY{o}{+} \PY{n}{d\PYZus{}H2O} \PY{o}{+} \PY{l+m+mi}{3}\PY{o}{*}\PY{n}{f\PYZus{}K2CO3}\PY{p}{,}
               \PY{l+s+s2}{\PYZdq{}}\PY{l+s+s2}{C}\PY{l+s+s2}{\PYZdq{}}\PY{p}{:} \PY{n}{c\PYZus{}CO2} \PY{o}{+} \PY{n}{f\PYZus{}K2CO3}\PY{p}{,} \PY{l+s+s2}{\PYZdq{}}\PY{l+s+s2}{H}\PY{l+s+s2}{\PYZdq{}}\PY{p}{:} \PY{l+m+mi}{2}\PY{o}{*}\PY{n}{d\PYZus{}H2O}\PY{p}{\PYZcb{}}

\PY{k}{for} \PY{n}{elem} \PY{o+ow}{in} \PY{p}{[}\PY{l+s+s2}{\PYZdq{}}\PY{l+s+s2}{K}\PY{l+s+s2}{\PYZdq{}}\PY{p}{,} \PY{l+s+s2}{\PYZdq{}}\PY{l+s+s2}{N}\PY{l+s+s2}{\PYZdq{}}\PY{p}{,} \PY{l+s+s2}{\PYZdq{}}\PY{l+s+s2}{O}\PY{l+s+s2}{\PYZdq{}}\PY{p}{,} \PY{l+s+s2}{\PYZdq{}}\PY{l+s+s2}{C}\PY{l+s+s2}{\PYZdq{}}\PY{p}{,} \PY{l+s+s2}{\PYZdq{}}\PY{l+s+s2}{H}\PY{l+s+s2}{\PYZdq{}}\PY{p}{]}\PY{p}{:}
    \PY{n}{l}\PY{p}{,} \PY{n}{r} \PY{o}{=} \PY{n}{atoms\PYZus{}left}\PY{p}{[}\PY{n}{elem}\PY{p}{]}\PY{p}{,} \PY{n}{atoms\PYZus{}right}\PY{p}{[}\PY{n}{elem}\PY{p}{]}
    \PY{n}{status} \PY{o}{=} \PY{l+s+s2}{\PYZdq{}}\PY{l+s+s2}{✓}\PY{l+s+s2}{\PYZdq{}} \PY{k}{if} \PY{n}{l} \PY{o}{==} \PY{n}{r} \PY{k}{else} \PY{l+s+s2}{\PYZdq{}}\PY{l+s+s2}{✗}\PY{l+s+s2}{\PYZdq{}}
    \PY{n+nb}{print}\PY{p}{(}\PY{l+s+sa}{f}\PY{l+s+s2}{\PYZdq{}}\PY{l+s+s2}{  }\PY{l+s+si}{\PYZob{}}\PY{n}{elem}\PY{l+s+si}{\PYZcb{}}\PY{l+s+s2}{: left=}\PY{l+s+si}{\PYZob{}}\PY{n}{l}\PY{l+s+si}{:}\PY{l+s+s2}{4d}\PY{l+s+si}{\PYZcb{}}\PY{l+s+s2}{  right=}\PY{l+s+si}{\PYZob{}}\PY{n}{r}\PY{l+s+si}{:}\PY{l+s+s2}{4d}\PY{l+s+si}{\PYZcb{}}\PY{l+s+s2}{  }\PY{l+s+si}{\PYZob{}}\PY{n}{status}\PY{l+s+si}{\PYZcb{}}\PY{l+s+s2}{\PYZdq{}}\PY{p}{)}
\end{Verbatim}
\end{tcolorbox}

    \begin{Verbatim}[commandchars=\\\{\}]

Balanced combustion equation:
  26 KNO₃  +  5 C₆H₁₄O₆  →  17 CO₂  +  35 H₂O  +  13 N₂  +  13 K₂CO₃

Atom-conservation check:
  K: left=  26  right=  26  ✓
  N: left=  26  right=  26  ✓
  O: left= 108  right= 108  ✓
  C: left=  30  right=  30  ✓
  H: left=  70  right=  70  ✓
    \end{Verbatim}

    \begin{center}\rule{0.5\linewidth}{0.5pt}\end{center}

\subsection{4. Stoichiometric mass ratio
derivation}\label{stoichiometric-mass-ratio-derivation}

\subsubsection{4.1 O₂-demand approach}\label{oux2082-demand-approach}

The O₂ demand of sorbitol (complete combustion) is:

\[
\text{C}_6\text{H}_{14}\text{O}_6 + 6\,\text{O}_2
\;\rightarrow\;
6\,\text{CO}_2 + 7\,\text{H}_2\text{O}
\]

Each mole of KNO₃ supplies 3 oxygen atoms = 1.5 mol O₂-equivalent, so
the stoichiometric KNO₃:sorbitol molar ratio is \(6 / 1.5 = 4\).

    \begin{tcolorbox}[breakable, size=fbox, boxrule=1pt, pad at break*=1mm,colback=cellbackground, colframe=cellborder]
\prompt{In}{incolor}{7}{\boxspacing}
\begin{Verbatim}[commandchars=\\\{\}]
\PY{c+c1}{\PYZsh{} ── Molar masses (g/mol) ──────────────────────────────────────────────────}
\PY{n}{M\PYZus{}KNO3}   \PY{o}{=} \PY{l+m+mf}{39.102} \PY{o}{+} \PY{l+m+mf}{14.007} \PY{o}{+} \PY{l+m+mi}{3} \PY{o}{*} \PY{l+m+mf}{15.999}          \PY{c+c1}{\PYZsh{} K + N + 3O}
\PY{n}{M\PYZus{}sorb}   \PY{o}{=} \PY{l+m+mi}{6}\PY{o}{*}\PY{l+m+mf}{12.011} \PY{o}{+} \PY{l+m+mi}{14}\PY{o}{*}\PY{l+m+mf}{1.008} \PY{o}{+} \PY{l+m+mi}{6}\PY{o}{*}\PY{l+m+mf}{15.999}         \PY{c+c1}{\PYZsh{} C6H14O6}

\PY{n+nb}{print}\PY{p}{(}\PY{l+s+sa}{f}\PY{l+s+s2}{\PYZdq{}}\PY{l+s+s2}{Molar mass KNO₃   : }\PY{l+s+si}{\PYZob{}}\PY{n}{M\PYZus{}KNO3}\PY{l+s+si}{:}\PY{l+s+s2}{.5f}\PY{l+s+si}{\PYZcb{}}\PY{l+s+s2}{ g/mol}\PY{l+s+s2}{\PYZdq{}}\PY{p}{)}
\PY{n+nb}{print}\PY{p}{(}\PY{l+s+sa}{f}\PY{l+s+s2}{\PYZdq{}}\PY{l+s+s2}{Molar mass sorbitol: }\PY{l+s+si}{\PYZob{}}\PY{n}{M\PYZus{}sorb}\PY{l+s+si}{:}\PY{l+s+s2}{.5f}\PY{l+s+si}{\PYZcb{}}\PY{l+s+s2}{ g/mol}\PY{l+s+s2}{\PYZdq{}}\PY{p}{)}

\PY{c+c1}{\PYZsh{} ── Molar ratio from balanced equation ───────────────────────────────────}
\PY{n}{molar\PYZus{}ratio\PYZus{}KNO3\PYZus{}to\PYZus{}sorb} \PY{o}{=} \PY{n}{a\PYZus{}KNO3} \PY{o}{/} \PY{n}{b\PYZus{}C6H14O6}
\PY{n+nb}{print}\PY{p}{(}\PY{l+s+sa}{f}\PY{l+s+s2}{\PYZdq{}}\PY{l+s+se}{\PYZbs{}n}\PY{l+s+s2}{Molar ratio KNO₃ : sorbitol  =  }\PY{l+s+si}{\PYZob{}}\PY{n}{a\PYZus{}KNO3}\PY{l+s+si}{\PYZcb{}}\PY{l+s+s2}{ : }\PY{l+s+si}{\PYZob{}}\PY{n}{b\PYZus{}C6H14O6}\PY{l+s+si}{\PYZcb{}}\PY{l+s+s2}{\PYZdq{}}
      \PY{l+s+sa}{f}\PY{l+s+s2}{\PYZdq{}}\PY{l+s+s2}{  =  }\PY{l+s+si}{\PYZob{}}\PY{n}{molar\PYZus{}ratio\PYZus{}KNO3\PYZus{}to\PYZus{}sorb}\PY{l+s+si}{:}\PY{l+s+s2}{.1f}\PY{l+s+si}{\PYZcb{}}\PY{l+s+s2}{ : 1}\PY{l+s+s2}{\PYZdq{}}\PY{p}{)}

\PY{c+c1}{\PYZsh{} ── Mass of KNO₃ needed per mol sorbitol ─────────────────────────────────}
\PY{n}{mass\PYZus{}KNO3\PYZus{}per\PYZus{}mol\PYZus{}sorb} \PY{o}{=} \PY{n}{molar\PYZus{}ratio\PYZus{}KNO3\PYZus{}to\PYZus{}sorb} \PY{o}{*} \PY{n}{M\PYZus{}KNO3}
\PY{n+nb}{print}\PY{p}{(}\PY{l+s+sa}{f}\PY{l+s+s2}{\PYZdq{}}\PY{l+s+s2}{Mass of KNO₃ per mol sorbitol: }\PY{l+s+si}{\PYZob{}}\PY{n}{mass\PYZus{}KNO3\PYZus{}per\PYZus{}mol\PYZus{}sorb}\PY{l+s+si}{:}\PY{l+s+s2}{.5f}\PY{l+s+si}{\PYZcb{}}\PY{l+s+s2}{ g}\PY{l+s+s2}{\PYZdq{}}\PY{p}{)}

\PY{c+c1}{\PYZsh{} ── Ideal mass fractions ──────────────────────────────────────────────────}
\PY{n}{total\PYZus{}mass} \PY{o}{=} \PY{n}{mass\PYZus{}KNO3\PYZus{}per\PYZus{}mol\PYZus{}sorb} \PY{o}{+} \PY{n}{M\PYZus{}sorb}
\PY{n}{wt\PYZus{}KNO3}    \PY{o}{=} \PY{n}{mass\PYZus{}KNO3\PYZus{}per\PYZus{}mol\PYZus{}sorb} \PY{o}{/} \PY{n}{total\PYZus{}mass}
\PY{n}{wt\PYZus{}sorb}    \PY{o}{=} \PY{n}{M\PYZus{}sorb} \PY{o}{/} \PY{n}{total\PYZus{}mass}

\PY{n+nb}{print}\PY{p}{(}\PY{l+s+sa}{f}\PY{l+s+s2}{\PYZdq{}}\PY{l+s+se}{\PYZbs{}n}\PY{l+s+s2}{Ideal stoichiometric mass fractions:}\PY{l+s+s2}{\PYZdq{}}\PY{p}{)}
\PY{n+nb}{print}\PY{p}{(}\PY{l+s+sa}{f}\PY{l+s+s2}{\PYZdq{}}\PY{l+s+s2}{  KNO₃    :  }\PY{l+s+si}{\PYZob{}}\PY{n}{wt\PYZus{}KNO3}\PY{o}{*}\PY{l+m+mi}{100}\PY{l+s+si}{:}\PY{l+s+s2}{.2f}\PY{l+s+si}{\PYZcb{}}\PY{l+s+s2}{ \PYZpc{}}\PY{l+s+s2}{\PYZdq{}}\PY{p}{)}
\PY{n+nb}{print}\PY{p}{(}\PY{l+s+sa}{f}\PY{l+s+s2}{\PYZdq{}}\PY{l+s+s2}{  Sorbitol:  }\PY{l+s+si}{\PYZob{}}\PY{n}{wt\PYZus{}sorb}\PY{o}{*}\PY{l+m+mi}{100}\PY{l+s+si}{:}\PY{l+s+s2}{.2f}\PY{l+s+si}{\PYZcb{}}\PY{l+s+s2}{ \PYZpc{}}\PY{l+s+s2}{\PYZdq{}}\PY{p}{)}

\PY{c+c1}{\PYZsh{} ── Mazza\PYZhy{}style O/F ratio (oxidiser / fuel mass ratio) ───────────────────}
\PY{n}{OF\PYZus{}ratio} \PY{o}{=} \PY{n}{mass\PYZus{}KNO3\PYZus{}per\PYZus{}mol\PYZus{}sorb} \PY{o}{/} \PY{n}{M\PYZus{}sorb}
\PY{n+nb}{print}\PY{p}{(}\PY{l+s+sa}{f}\PY{l+s+s2}{\PYZdq{}}\PY{l+s+se}{\PYZbs{}n}\PY{l+s+s2}{O/F mass ratio (KNO₃ / sorbitol): }\PY{l+s+si}{\PYZob{}}\PY{n}{OF\PYZus{}ratio}\PY{l+s+si}{:}\PY{l+s+s2}{.4f}\PY{l+s+si}{\PYZcb{}}\PY{l+s+s2}{\PYZdq{}}\PY{p}{)}
\end{Verbatim}
\end{tcolorbox}

    \begin{Verbatim}[commandchars=\\\{\}]
Molar mass KNO₃   : 101.10600 g/mol
Molar mass sorbitol: 182.17200 g/mol

Molar ratio KNO₃ : sorbitol  =  26 : 5  =  5.2 : 1
Mass of KNO₃ per mol sorbitol: 525.75120 g

Ideal stoichiometric mass fractions:
  KNO₃    :  74.27 \%
  Sorbitol:  25.73 \%

O/F mass ratio (KNO₃ / sorbitol): 2.8860
    \end{Verbatim}

    \subsubsection{4.2 Summary}\label{summary}

{\def\LTcaptype{none} % do not increment counter
\begin{longtable}[]{@{}ll@{}}
\toprule\noalign{}
Quantity & Value \\
\midrule\noalign{}
\endhead
\bottomrule\noalign{}
\endlastfoot
Stoichiometric molar ratio \(n_\text{KNO₃} : n_\text{sorb}\) & 26 : 5 \\
Stoichiometric O/F mass ratio & ≈ 2.22 \\
\textbf{KNO₃ mass fraction} & \textbf{≈ 68.9 \%} \\
\textbf{Sorbitol mass fraction} & \textbf{≈ 31.1 \%} \\
\end{longtable}
}

These are the \emph{stoichiometric} values --- full combustion of
sorbitol with the oxygen supplied by KNO₃.

    \begin{center}\rule{0.5\linewidth}{0.5pt}\end{center}

\subsection{5. Why the operational mix is slightly fuel-rich (65\% :
35\%)}\label{why-the-operational-mix-is-slightly-fuel-rich-65-35}

It has been established experimentally and confirmed by thermochemical
codes (e.g.~NASA-CEA) that a slightly \textbf{fuel-rich} blend ---
typically 65\% KNO₃ / 35\% sorbitol --- outperforms the stoichiometric
point in several respects:

    \begin{tcolorbox}[breakable, size=fbox, boxrule=1pt, pad at break*=1mm,colback=cellbackground, colframe=cellborder]
\prompt{In}{incolor}{8}{\boxspacing}
\begin{Verbatim}[commandchars=\\\{\}]
\PY{n}{operational} \PY{o}{=} \PY{p}{\PYZob{}}
    \PY{l+s+s2}{\PYZdq{}}\PY{l+s+s2}{KNO₃}\PY{l+s+s2}{\PYZdq{}}    \PY{p}{:} \PY{l+m+mf}{0.65}\PY{p}{,}
    \PY{l+s+s2}{\PYZdq{}}\PY{l+s+s2}{Sorbitol}\PY{l+s+s2}{\PYZdq{}} \PY{p}{:} \PY{l+m+mf}{0.35}\PY{p}{,}
\PY{p}{\PYZcb{}}

\PY{n+nb}{print}\PY{p}{(}\PY{l+s+s2}{\PYZdq{}}\PY{l+s+s2}{Operational KNSB mix (mass fractions):}\PY{l+s+s2}{\PYZdq{}}\PY{p}{)}
\PY{k}{for} \PY{n}{species}\PY{p}{,} \PY{n}{frac} \PY{o+ow}{in} \PY{n}{operational}\PY{o}{.}\PY{n}{items}\PY{p}{(}\PY{p}{)}\PY{p}{:}
    \PY{n+nb}{print}\PY{p}{(}\PY{l+s+sa}{f}\PY{l+s+s2}{\PYZdq{}}\PY{l+s+s2}{  }\PY{l+s+si}{\PYZob{}}\PY{n}{species}\PY{l+s+si}{:}\PY{l+s+s2}{\PYZlt{}12}\PY{l+s+si}{\PYZcb{}}\PY{l+s+s2}{: }\PY{l+s+si}{\PYZob{}}\PY{n}{frac}\PY{o}{*}\PY{l+m+mi}{100}\PY{l+s+si}{:}\PY{l+s+s2}{.1f}\PY{l+s+si}{\PYZcb{}}\PY{l+s+s2}{ \PYZpc{}}\PY{l+s+s2}{\PYZdq{}}\PY{p}{)}

\PY{n+nb}{print}\PY{p}{(}\PY{l+s+sa}{f}\PY{l+s+s2}{\PYZdq{}}\PY{l+s+se}{\PYZbs{}n}\PY{l+s+s2}{Deviation from stoichiometric KNO₃ fraction:  }\PY{l+s+s2}{\PYZdq{}}
      \PY{l+s+sa}{f}\PY{l+s+s2}{\PYZdq{}}\PY{l+s+si}{\PYZob{}}\PY{p}{(}\PY{n}{operational}\PY{p}{[}\PY{l+s+s1}{\PYZsq{}}\PY{l+s+s1}{KNO₃}\PY{l+s+s1}{\PYZsq{}}\PY{p}{]}\PY{+w}{ }\PY{o}{\PYZhy{}}\PY{+w}{ }\PY{n}{wt\PYZus{}KNO3}\PY{p}{)}\PY{o}{*}\PY{l+m+mi}{100}\PY{l+s+si}{:}\PY{l+s+s2}{+.2f}\PY{l+s+si}{\PYZcb{}}\PY{l+s+s2}{ pp}\PY{l+s+s2}{\PYZdq{}}\PY{p}{)}
\end{Verbatim}
\end{tcolorbox}

    \begin{Verbatim}[commandchars=\\\{\}]
Operational KNSB mix (mass fractions):
  KNO₃        : 65.0 \%
  Sorbitol    : 35.0 \%

Deviation from stoichiometric KNO₃ fraction:  -9.27 pp
    \end{Verbatim}

    The excess sorbitol produces the following thermochemical advantages:

{\def\LTcaptype{none} % do not increment counter
\begin{longtable}[]{@{}
  >{\raggedright\arraybackslash}p{(\linewidth - 4\tabcolsep) * \real{0.3333}}
  >{\raggedright\arraybackslash}p{(\linewidth - 4\tabcolsep) * \real{0.3333}}
  >{\raggedright\arraybackslash}p{(\linewidth - 4\tabcolsep) * \real{0.3333}}@{}}
\toprule\noalign{}
\begin{minipage}[b]{\linewidth}\raggedright
Effect
\end{minipage} & \begin{minipage}[b]{\linewidth}\raggedright
Mechanism
\end{minipage} & \begin{minipage}[b]{\linewidth}\raggedright
Propulsion benefit
\end{minipage} \\
\midrule\noalign{}
\endhead
\bottomrule\noalign{}
\endlastfoot
Sub-adiabatic flame temperature & Incomplete oxidation absorbs enthalpy
& Reduced casing thermal load \\
More gas-phase products (CO, H₂) & Incomplete combustion raises moles of
gas & Increased \(I_{sp}\) \\
Lower \(\text{K}_2\text{CO}_3\) yield & Oxygen limited below
stoichiometric & Reduced nozzle slag / clogging \\
Stickier, more plastic slurry & Higher binder (sugar) loading & Better
castability and grain integrity \\
\end{longtable}
}

The trade-off is a marginal reduction in total energy release; the above
benefits are judged to outweigh this for hobby- and small
experimental-rocket applications (Nakka 2025).

    \begin{center}\rule{0.5\linewidth}{0.5pt}\end{center}

\subsection{6. Conclusion}\label{conclusion}

Starting from the Mazza nullspace method (MATLAB, Chapter 2), we
produced a clean Python/SymPy drop-in replacement that:

\begin{itemize}
\tightlist
\item
  parses chemical formulas directly from symbolic expressions --- no
  manual matrix entry;
\item
  builds the signed balance matrix \(\mathbf{A}\) algorithmically;
\item
  computes the exact integer nullspace using
  \texttt{Matrix.nullspace()};
\item
  detects underdetermined systems (nullspace dim \textgreater{} 1) and
  reports them as warnings.
\end{itemize}

Applied to KNSB, the method confirmed that the 9-species product slate
is underdetermined and returned the \textbf{unique balanced equation}
for the simplified 6-species system:

\[
\boxed{
26\,\text{KNO}_3 + 5\,\text{C}_6\text{H}_{14}\text{O}_6
\;\longrightarrow\;
17\,\text{CO}_2 + 35\,\text{H}_2\text{O} + 13\,\text{N}_2 + 13\,\text{K}_2\text{CO}_3
}
\]

The stoichiometric mass fraction of KNO₃ is \textbf{≈ 68.9\%}, and the
operational blend (65\% KNO₃ / 35\% sorbitol) is chosen to sit
\textasciitilde4 percentage points fuel-rich of the stoichiometric point
for thermal, slag-minimisation, and castability reasons.

    \begin{center}\rule{0.5\linewidth}{0.5pt}\end{center}

\subsection{References}\label{references}

\begin{enumerate}
\def\labelenumi{\arabic{enumi}.}
\tightlist
\item
  D. Mazza and E. Canuto, \emph{Fundamental Chemistry with MATLAB},
  Elsevier, 2022.
\item
  R. Nakka, ``KNSB Propellant,'' \emph{Richard Nakka's Experimental
  Rocketry Web Site}, 2025. {[}Online{]}. Available:
  https://www.nakka-rocketry.net/sorb.html
\item
  J. Bonnie, J. Zehe, and S. Gordon, \emph{NASA Glenn Coefficients for
  Calculating Thermodynamic Properties of Individual Species},
  NASA/TP-2002-211556, Glenn Research Center, Cleveland, 2002.
\item
  W. Givens, ``Parametric solution of linear homogeneous Diophantine
  equations,'' \emph{Bull. Am. Math. Soc.}, vol.~53, no. 8,
  pp.~780--783, 1947.
\end{enumerate}


    % Add a bibliography block to the postdoc
    
    
    
\end{document}
